% !TEX root = ../nc-cubic-n-fold.tex
\section{残余比較のための Poisson 補正}
\label{sec:nfm-poisson-correction}

比較障害の後に残る問題は、残余圏が復元する三次と同じ三次を選ぶ
Poisson 対を、一般次元の族として構成することである。
まず全 O-同値に降下する代数的補正を与え、次に固定 formality による
全次数の形式構成と比較する。最後に分離型の族では、固定入力の下で
代数的な補正後三次への全次数の残余比較を示す。
三つの構成の適用範囲を区別し、一般混合係数での全次数の一致を仮定しない。

\subsection{Hessian と trace による代数的な補正}
\label{sec:nfm-algebraic-correction}

有限 jet の障害は、周囲 Poisson tensor を変えることで取り除ける。
ここではその二次 jet を持つ、座標に依存しない代数的な処方を構成する。
この処方の全次数が残余復元と一致するとは主張しない。

$R$ を可換 $\mathbf C$-代数、$V$ を階数 $N=n+2$ の有限局所自由加群、
$W=V^*$ とする。分解可能 Poisson 対を局所的に $(H,P)$ と表し、
$K=\operatorname{supp}(P)\subset W$、$U=\operatorname{ann}(K)\subset V$ と置く。
$H$ の定める三次超曲面と、その $\mathbb P(K)$ への制限がともに滑らかとする。
$S=\operatorname{Sym}U$ および
\[
 T_{H,K}=\operatorname{Sym}V/(\partial_kH:k\in K),\qquad
 \tau_{H,K}=\tfrac18\operatorname{Tr}_{T_{H,K}/S}
\]
と置く。ここでも trace には商への写像を合成する。
active な勾配の最高次部分が三変数の正則列をなすので、
$T_{H,K}$ は $S$ 上階数八の有限局所自由代数である。
この性質と trace は任意の基底変更と両立する。

$K$ 方向の二階微分は
$D_K^2H\in\operatorname{Sym}^2K^*\otimes V$ であり、その行列式は
$(\det K^*)^{\otimes2}\otimes\operatorname{Sym}^3V$ に属する。
従って次の三次式は $K$ の座標と体積形式の選択に依存しない。
\begin{equation}\label{eq:nfm-intrinsic-hessian-correction}
 \mathcal D_P(H)=\langle\det(D_K^2H),P\otimes P\rangle,
 \qquad
 \mathcal C_P(H)=\frac1{24}\mathcal D_P(H)
       +\frac1{12}\tau_{H,K}(\mathcal D_P(H)).
\end{equation}
例えば $P=t\,\partial_x\wedge\partial_y\wedge\partial_z$ なら
$\mathcal D_P(H)=t^2\det\operatorname{Hess}_{x,y,z}(H)$ である。
trace 項は $\operatorname{Sym}^3U$ に属する。

\begin{theorem}[良い Poisson 対の代数的 Hessian 補正]
\label{thm:nfm-algebraic-hessian-correction}
active な平面三次も滑らかな $\mathfrak P_n^{\mathrm{dec}}$ の開を
$\mathfrak P_n^{\mathrm{ell}}$ と書く。
次の処方は、出力の三次と active な平面三次がともに滑らかな非空開上で、
代数的な射
\[
 \Phi:\mathfrak P_n^{\mathrm{ell}}\dashrightarrow
             \mathfrak P_n^{\mathrm{ell}},\qquad
 (H,P)\longmapsto(H^\sharp,P),\quad H^\sharp=H+\mathcal C_P(H)
\]
を定める。これは任意の基底変更、方程式の生成元変更、
全ての射影座標同型と両立する。
補正後の強い lift は厳密に
\begin{equation}\label{eq:nfm-corrected-algebraic-tensor}
 \Pi^\sharp=\iota_{dH^\sharp}P
 =\Pi+\frac1{24}\iota_{d\mathcal D_P(H)}P
\end{equation}
であり、同じ $K$ と $U$ を持つ Poisson tensor である。

この有理射は支配的で、非空開への制限は \emph{\'etale} である。
さらに定理\ref{thm:nfm-correction-degree-nine}で、始域と終域を適切に縮めると
有限 \emph{\'etale} で階数 $9$ となることを示す。
従って定理\ref{thm:poisson-good-open-isomorphism}を介し、
$\mathfrak M_n^{\mathrm{sm}}$ の非空開から Poisson 対のモジュライへの射を与える。
必要なら定義域をさらに縮め、出力も $\mathfrak P_n^{\mathrm{good}}$ に入るようにできる。
ただしこのことは、元の非可換圏とその出力の Kuznetsov 成分の
全次数の同値を意味しない。
\end{theorem}

\begin{proof}
三次方程式を $H\mapsto cH$ とすると、強い lift の表示は
$P\mapsto c^{-1}P$ となる。
Hessian 行列式は $c^3$ 倍、$P^{\otimes2}$ は $c^{-2}$ 倍になるから
\[
 \mathcal D_{c^{-1}P}(cH)=c\mathcal D_P(H).
\]
active な勾配イデアルは変わらず、trace も変わらない。
従って $\mathcal C_{c^{-1}P}(cH)=c\mathcal C_P(H)$ であり、
出力 $H^\sharp\otimes P$ は方程式の生成元に依存しない。
Hessian の縮約と有限局所自由代数の trace は線型座標変換に共変である。
特に active 変数に中心一次式を加える frame 変更も含み、
$V\twoheadrightarrow K^*$ の分裂を選ぶ必要はない。
次数スカラーの重みも相殺するので、射影群の商 stack へ降下する。
有限局所自由性と trace の基底変更により、非平坦・非被約の場合も含めて
任意の基底変更と両立する。

局所的に $K$ の座標を選ぶと、$\Pi^\sharp$ は三変数の Jacobian 括弧で、
残りの変数は中心である。従って Jacobi 恒等式は多項式恒等式として成立し、
$H^\sharp$ はその Casimir である。
trace 項の微分は $K$ 方向で零なので、式
\eqref{eq:nfm-corrected-algebraic-tensor}を得る。
出力が滑らかなら、系\ref{cor:poisson-support-recovery}により一次中心は
ちょうど $U$ である。

定義域の非空性と支配性を一様に確認する。
滑らかな三次と active な平面三次の開
\[
 B\subset\mathbb P(\operatorname{Sym}^3V)\times\operatorname{Gr}(3,W)
\]
上で、$\mathcal L_H\otimes\det\mathcal K$ の全空間 $\overline Z$ を取り、
零切断も許す。ただし零切断では $K$ を忘れない。
式\eqref{eq:nfm-intrinsic-hessian-correction}は $P=0$ でも定義でき、
補正は線束方向の二次以上である。
したがって $\Phi$ は零切断を含む開で定義され、零切断を各点ごとに固定し、
そこでの微分は恒等である。
$\overline Z$ は滑らかなので、$\Phi$ は零切断の近傍で \emph{\'etale} となる。
この近傍の零切断を除いた部分は非空で、像は開かつ稠密である。
よって元の非零部分上の有理射は支配的であり、非空開で \emph{\'etale} である。
同じ群による同変商に移すと、この性質は stack 上にも降りる。

$\mathfrak P_n^{\mathrm{good}}$ は既約な
$\mathfrak P_n^{\mathrm{ell}}$ の非空開である。
支配性により、その逆像と定義域中の
$\mathfrak P_n^{\mathrm{good}}$ の交わりも非空開となる。
そこに既存の O-marked moduli との同型を適用すればよい。
\end{proof}

\begin{corollary}[flag を残した零切断の形式近傍]
\label{cor:nfm-correction-formal-automorphism}
上の $\overline Z$ の零切断に沿う形式完備化上で、$\Phi$ は
恒等を剰余に持つ形式自己同型である。
同じことは射影群の商にも成立する。
\end{corollary}
\begin{proof}
零切断を固定する \emph{\'etale} 射に対し、零切断上の恒等写像の
持上げは各冪零厚みで一意である。
これらを順次持ち上げると、任意の有限厚み上で相互に逆な写像を得る。
その逆極限が形式逆写像を与える。一意性により群作用とも両立する。
これは明示的な補正 $\Phi$ 自身の形式可逆性であり、
formality の中心写像との全次数の一致ではない。
\end{proof}

\begin{proposition}[楕円点曲線との厳密な整合]
\label{prop:nfm-correction-point-cubic}
良い一ブロック型の ambient 二次代数 $B_\Pi$ を考える。
$g=H|_{\mathbb P(K)}$ とすると、補正後の制限は
\[
 H^\sharp|_{\mathbb P(K)}
 =g+\frac1{24}\mathcal D_P(g)
\]
であり、これは $B_\Pi$ の active な楕円点曲線の方程式である。
従って補正射は、回復した support 平面の中で、
量子側の楕円点曲線と同じ平面三次を選ぶ。
この主張は点曲線の shift を保持した対応全体や、
定常点成分の同定を含まない。
\end{proposition}

三生成元の二次 Calabi--Yau 代数の点曲線を Hessian によって記述することは、
Mori--Smith\cite[Theorems~1.2, 7.2]{MoriSmith}の既知の結果である。
以下の証明は本稿の係数規約を確認するもので、その点曲線公式自体の新規性を主張しない。
本節ではこれを混合係数の trace と組み合わせ、
全次元の Poisson 対および O-marked moduli 上の代数的補正に用いる。

\begin{proof}
trace 項は純中心三次なので、$\mathbb P(K)$ への制限で零になる。
局所座標で $P=t\,\partial_x\wedge\partial_y\wedge\partial_z$ と書き、
$p=(x,y,z)^{\mathsf t}$ とする。二次関係の双線型化は
\[
 \left(C(p)-\frac t2\operatorname{Hess}(g)(p)\right)q=0,
 \qquad C(p)q=p\times q
\]
である。対称行列 $M$ に対する三次恒等式
\[
 \det(C(p)+sM)=s\,p^{\mathsf t}Mp+s^3\det M
\]
と $p^{\mathsf t}\operatorname{Hess}(g)(p)p=6g(p)$ より、
その行列式は
\[
 -3t\left(g+\frac{t^2}{24}\det\operatorname{Hess}(g)\right)
\]
である。非零の $P$ を持つ相対チャートでは $t$ は単元に取れる。
行列式恒等式は、$\operatorname{adj}C(p)=pp^{\mathsf t}$ と
転置による $s$ の偶数次係数の消滅からも直接分かる。
この三次が滑らかな開では、係数行列の階数は曲線上の全ての幾何点で二である。
階数が一以下なら全ての余因子が消え、行列式の全偏微分も消えるので、
滑らかさに反する。さらに $Mq=0$ に $\operatorname{adj}(M)$ を掛けると
$\det(M)q=0$ となり、射影座標の局所的な可逆成分から incidence は
行列式曲線をスキームとして経由する。
可逆な二次小行列を使えば、任意基底上で核は直線部分束となる。
従って核は任意基底変更と両立する直線部分束であり、
双線型関係の incidence はこの曲線からの射のグラフとなる。
左右の変数を交換すると $t$ の符号が変わるだけなので、
右射影の三次行列式も同じ曲線を与え、その核も直線部分束である。
従って両射影は同型で、incidence は曲線の自己同型のグラフである。
これにより三次が楕円点曲線であることも任意基底上で確認できる。
$t=0$ で得られる三次はこの族の極限であり、
可換多項式環の点スキームそのものと同一視しない。
\end{proof}

\begin{corollary}[残余復元の二次 jet を実現する代数的処方]
\label{cor:nfm-corrected-residual-jet}
命題\ref{prop:nfm-relative-residual-second-jet}の規約と入力を使う。
Hessian 正規化中心元の trace--Hessian 出力に $\Phi$ を適用すると、
\[
 H^\sharp
 =F+\frac{h^2}{24}\mathscr H_F+h^4\tau_F(\Psi_\lambda(F)),
 \qquad H^\sharp\equiv W^{\mathrm{res}}\pmod{h^4}.
\]
その強い lift は
\[
 \Pi^\sharp=h\Pi_F+\frac{h^3}{24}\Pi_{\mathscr H_F}
\]
である。従って前節の二次 Jacobi 障害を、この代数的な補正により除去できる。
\end{corollary}
\begin{proof}
trace--Hessian 出力と $F$ の差は純中心三次なので、
active な勾配、Hessian、および $\tau_F$ は変わらない。
その強い lift の表示は $P=hP_0$ である。
従って $\mathcal D_{hP_0}(H^{\mathrm{TH}})=h^2\mathscr H_F$ となり、
二つの trace 項が相殺して表示した式を得る。
tensor の式は純中心項の縮約が零であることから従う。
残余二次 jet と以下の系\ref{cor:nfm-mixed-even-jet}の偶奇性により、
合同の剰余は $h^4$ まで改善する。
この計算は $h\ne0$ の対に対する代数的処方を当該 trait に沿って延長したもので、
可換特殊点全体の O-marked moduli に非零 Poisson tensor を付与するものではない。
\end{proof}

\begin{remark}[残る比較問題]
$\Phi$ は特定の formality の全次数の逆写像として定義したものではない。
ここで示したのは、任意基底の代数的構成、全 O-同値への降下、
楕円点曲線との一致、および指定中心元に対する残余三次の二次 jet との一致である。
実際の formal Poisson tensor は scalar formal unit による規格化も持ち得る。
従ってその tensor との比較には、potential と support の比較に加えて尺度を追跡する。
全次数での $\operatorname{Ku}(V(H^\sharp))$ との同値、
またはこの処方の一意性は、二次 jet の一致からは従わない。
\end{remark}

\subsection{一般混合係数における補正射の次数}
\label{sec:nfm-correction-degree}

Hessian 補正の零切断に沿う形式逆写像は、大域の有理逆写像ではない。
この相違は分離型の族に限らず、任意次元の一般混合係数で定量化できる。
以下の次数は、同じ射影座標群による商 stack の間の表現可能射の次数である。
粗モジュライの関数体の次数を、慣性群の比較なしに言い換えてはいない。
Hesse pencil と同変性による還元は、射影的な純 Hessian 写像の次数三を扱う
Ciliberto--Ottaviani\cite[Section~5, Theorem~5.7]{NFMCilibertoOttaviani}にも現れる。
本稿で数えるのは $f\mapsto f+\det\operatorname{Hess}(f)/24$ とその混合拡張であり、
純 Hessian 写像の二回合成の次数九とも異なる。

\begin{theorem}[補正射の一般次数]
\label{thm:nfm-correction-degree-nine}
任意の $n\geq3$ に対し、代数的 Hessian--trace 補正
\[
 \Phi:\mathfrak P_n^{\mathrm{ell}}\dashrightarrow
       \mathfrak P_n^{\mathrm{ell}}
\]
の一般次数は $9$ である。
より正確には、始域と終域に非空開部分 stack を取ると、$\Phi$ は
表現可能な有限 \emph{\'etale} 射で、階数は $9$ となる。
両方の開を $\mathfrak P_n^{\mathrm{good}}$ に含めてよい。
従って既存の $\mathfrak M_n^{\mathrm{sm}}\simeq
\mathfrak P_n^{\mathrm{good}}$ を介しても、同じ次数の射を得る。
特に、この補正射自身は有理同型ではない。
\end{theorem}

\begin{proof}
まず support $K$ を固定し、その体積三重ベクトル $P_0$ を選ぶ。
局所表示 $P=sP_0$ に対し、方程式の生成元変更によって
$(H,sP_0)$ を $(F,P_0)$、$F=sH$ と同一視する。
この表示における補正は
\begin{equation}\label{eq:nfm-normalized-correction-map}
 F\longmapsto F+\frac1{24}\Delta_F
                  +\frac1{12}\tau_F(\Delta_F),
 \qquad \Delta_F=\det\operatorname{Hess}_K(F).
\end{equation}
trace は $F$ の純中心部分に依存しない。

\smallskip\noindent
\textbf{active な三次の次数。}
三次元の体積座標 $(x,y,z)$ における多項式写像
\[
 \phi_3:\operatorname{Sym}^3\mathbf C^3{}^*
   \longrightarrow\operatorname{Sym}^3\mathbf C^3{}^*,
 \qquad f\longmapsto f+\tfrac1{24}\det\operatorname{Hess}(f)
\]
を考える。$g_{p,q}=p(x^3+y^3+z^3)/3+qxyz$ とおくと
\[
 \phi_3(g_{p,q})=g_{\psi_1(p,q),\psi_2(p,q)},\qquad
 \psi(p,q)=\left(p(1-q^2/4),\ q+p^3/3+q^3/12\right).
\]
この $\psi$ は、射影平面の射
\begin{equation}\label{eq:nfm-projective-correction-chart}
 [p:q:t]\longmapsto
 [3p(4t^2-q^2):12qt^2+4p^3+q^3:12t^3]
\end{equation}
へ拡張する。$t\ne0$ では最後の成分が零でなく、$t=0$ では
$pq^2=4p^3+q^3=0$ の非零解はないので基点はない。
この射は $\mathcal O(1)$ を $\mathcal O(3)$ に引き戻す。
従って有限かつ次数 $3^2=9$ である。
無限遠直線の逆像は同じ無限遠直線なので、$\psi$ 自身も有限であり、
滑らかな二次元の始域と終域の間の有限射として平坦、階数 $9$ となる。

ここで
\[
 \vartheta:\operatorname{SL}_3\times\mathbb A^2
    \longrightarrow\operatorname{Sym}^3\mathbf C^3{}^*,
 \qquad (A,p,q)\longmapsto A\cdot g_{p,q}
\]
は支配的かつ一般に有限である。
実際、$(1,1,0)$ における微分は、$\mathfrak{sl}_3$ の非対角成分から
六つの $x_i^2x_j$、対角成分から二つの $x_i^3$ の差、$p,q$ から
$x^3+y^3+z^3$ と $xyz$ を与える。これらは十次元の三次式空間の基底である。
Hessian は $\operatorname{SL}_3$ 同変なので
\[
 \vartheta\circ(1\times\psi)=\phi_3\circ\vartheta.
\]
一般次数の乗法性により
\[
 (\deg\vartheta)\,9=(\deg\phi_3)(\deg\vartheta),
 \qquad \deg\phi_3=9.
\]
この論証は Hesse frame の有限同定の個数を仮定していない。

\smallskip\noindent
\textbf{混合係数の復元。}
補空間を一時的に選び、active 次数で
\[
 F=f_3+f_2+f_1+f_0,\qquad
 f_i\in\operatorname{Sym}^iK^*\otimes
                    \operatorname{Sym}^{3-i}U
\]
と書く。$\operatorname{Hess}_K(F)$ は $f_3+f_2$ だけに依存する。
従って $\Delta_F=\delta_3+\delta_2+\delta_1+\delta_0$ と分けると
\[
 \delta_3=\det\operatorname{Hess}(f_3),\qquad
 \delta_2=\operatorname{tr}\bigl(
   \operatorname{adj}(\operatorname{Hess}(f_3))
                       \operatorname{Hess}_K(f_2)\bigr),
\]
また $\delta_1,\delta_0$ は $f_3,f_2$ だけから決まる。
式\eqref{eq:nfm-normalized-correction-map}は次の三角形の形になる。
\begin{align*}
 f'_3&=\phi_3(f_3),\\
 f'_2&=(L_{f_3}\otimes1_U)f_2,\\
 f'_1&=f_1+\delta_1/24,\\
 f'_0&=f_0+\delta_0/24+\tau_F(\Delta_F)/12,
\end{align*}
ここで
\[
 L_f(b)=b+\tfrac1{24}\operatorname{tr}\bigl(
   \operatorname{adj}(\operatorname{Hess}(f))\operatorname{Hess}(b)\bigr)
 \quad (b\in\operatorname{Sym}^2K^*).
\]
$L_f$ は一般の $f$ で可逆である。
例えば $f=g_{p,q}$ の場合、順序 $(x^2,yz)$、$(y^2,zx)$、$(z^2,xy)$
では同じ二次正方行列の三つの直和になり、各行列は
\[
 \begin{pmatrix}
  1-q^2/12&-pq/6\\ p^2/3&1+q^2/12
 \end{pmatrix}.
 \qquad
 \det L_{g_{p,q}}
   =\left(1+\frac{8p^3q-q^4}{144}\right)^3.
\]
特に $g_{1,0}$ で行列式は $1$ である。
一般の $f'_3$ には九つの $f_3$ があり、それらすべてで $L_{f_3}$ を
可逆にするよう終域の開を縮めてよい。
各 $f_3$ を選ぶと、上式から $f_2$、次に $f_1$、最後に $f_0$ が一意に決まる。
最後の trace は $f_0$ に依存しないため、そこで新たな方程式は発生しない。
よって任意の $\dim U=n-1$ に対して一般次数は同じく $9$ である。

\smallskip\noindent
\textbf{開の選択と商。}
ここまでの計算は有限型既約な係数空間の一般点の計算である。
一般に有限な写像では、始域の真の閉集合は終域を支配できない。
従って full cubic の特異点条件や good 開の補集合を除いても、
終域を縮めれば九つの一般分枝はすべて残る。
active 三次の滑らかさにより、trace を定める次数八の代数は有限局所自由で、
trace は正則である。全三次の滑らかさ、補正後の滑らかさ、
および既存の good 条件もそれぞれ非空開である。
一般に有限な写像では、取り除いた真の閉集合の像の閉包は終域を覆わない。
したがって有限性の開を取り、分岐像およびこれらの除外集合の像を取り除くと、
有限 \emph{\'etale} 階数 $9$ の射が得られる。
階数は一般の混合係数から求めており、特殊な分離型部分空間の次数を
そのまま全空間へ移したのではない。

表示 $(H,P)$ の方程式スカラーを商にした係数空間は、
$\operatorname{Gr}(3,W)$ 上の
$\operatorname{Sym}^3V\otimes\det\mathcal K$ の非零部分の開と見られる。
実際、その点は $H\otimes P$ である。
$\operatorname{PGL}(W)$ は support $K$ に推移的に作用し、
$K$ を固定して $P_0$ を選ぶと上の $F\otimes P_0$ 表示を得る。
$K$ の安定化群の作用は、体積の変化を $F$ の倍率へ吸収した作用であり、
始域と終域で同じである。従って support を固定した計算の次数 $9$ は、
support と体積 frame を動かした係数空間でも変わらない。
構成は同じ射影座標群 $G$ に関して同変であり、
有限 \emph{\'etale} 性を持つ終域の開を $G$ 不変に選べる。
一般に $G$ 同変な $X\to Y$ に対して
\[
 [X/G]\mathbin{\times}_{[Y/G]}Y\simeq X
\]
である。従って $[X/G]\to[Y/G]$ は表現可能であり、
有限 \emph{\'etale} 性と階数 $9$ はこの忠実平坦被覆から降下する。
以上により主張を得る。
\end{proof}

\begin{remark}[Hesse pencil と楕円曲線の乗法写像]
\label{rem:nfm-degree-not-multiplication}
active な入力と出力は同じ Hesse pencil に属するが、一般には同じ楕円曲線ではない。
例えば $q=0$、$p\ne0$ の入力曲線は Fermat 曲線である一方、出力の
Hesse 係数は $q'/p'=p^2/3$ であり、その $j$ 不変量は非定数である。
従って上の次数 $9$ を、入力の楕円曲線上の $[3]$ と同一視できない。
Hesse pencil とその $j$ 関数については
\cite[Section~2, Lemma~2.1, equations~(4)--(5)]{NFMArtebaniDolgachev}を参照。
Sklyanin 中心三次を別の superpotential と見たときに現れる
$(E,p)\mapsto(E,[-2]p)$\cite[Section~2]{NFMDLL}とも異なる。
後者は同じ $E$ を保つ。
今回の次数定理は、形式逆量子化の代数化や残余圏の全次数比較を主張していない。
\end{remark}

\subsection{残余復元と両立する形式 Poisson 対}
\label{sec:nfm-corrected-formal-poisson}

前節の障害は、量子代数の二次関係から直読する tensor と、
formality の逆対応が与える tensor を区別することで解消される。
本節では harmonic-angle formality を一つ固定し、形式パラメータ $h$ と
次数一生成元を保つ模型を用いる。
次数付き二次代数と二次 Poisson tensor の形式対応そのものは
Kontsevich\cite[Section~2.3, Proposition~6]{NFMKontsevichAlg}の既知結果である。
ここで明らかにするのは、一ブロックの場合の中心一次空間の保存、
残余三次による全次数の tensor の決定、および前節の障害を打ち消す補正である。
構成は形式的なものであり、非零の複素数 $h$ への収束・代数化とは区別する。

\subsubsection{次数一を固定する模型と有限係数の再帰}

$R$ を複素代数、$V_R$ を階数 $N$ の有限自由加群とし、
$S_R=\operatorname{Sym}_R V_R$ と置く。
次数一生成元を $z_1,\ldots,z_N$ とする。
次数別に $h$ 進完備な二次代数
\[
 B_h=T_R(V_R)[[h]]\Big/
 \left(z_i z_j-z_j z_i-h\operatorname{sym}(q_{ij})\right),
 \qquad q_{ij}\in\operatorname{Sym}^2_R V_R[[h]],\quad q_{ji}=-q_{ij},
\]
を考える。完備化は各内部次数で行い、最後にその直和を取る。
PBW 条件とは、この各内部次数が $S_R[[h]]$ と同じ有限局所自由加群になることをいう。
以下で存在を用いるのは、固定入力論文Bの
補題\texttt{cft:lem:model}が適用できる完備模型、
または後述の被約有限型の普遍係数模型からの基底変更である。

\begin{lemma}[生成元を固定する正規化の一意性]
\label{lem:nfm-normalized-star-model}
固定した affine-equivariant な Kontsevich 星積\cite[Section~6.3]{NFMKontsevich}について、
$h$ を固定し、全ての次数一生成元を固定する厳密同型
\[
 \phi:B_h\xrightarrow{\sim}(S_R[[h]],*_{\pi_h})
\]
が存在するなら、二次係数の formal Poisson tensor $\pi_h$ と $\phi$ は一意である。
もとの formality 同型が次数を保ち剰余で恒等なら、
形式的な線型座標変換によりこの正規化を行える。
従って正規化された対応は、係数に依存する可逆な線型座標変換にも共変である。
\end{lemma}

\begin{proof}
もとの同型の次数一部分は、剰余で恒等な可逆行列である。
その逆による線型座標変換を標的へ合成する。
グラフ写像の線型共変性により、標的は変換した tensor の同じ Kontsevich 星積であり、
全生成元が固定される。

二つの正規化された tensor が初めて $h^m$ で異なり、
$\pi'_h-\pi_h=h^mD\pmod{h^{m+1}}$ とする。
二頂点以上のグラフにこの差を入れると、他の tensor 入力が $h$ の倍数なので
$h^{m+1}$ の倍数になる。従って星積の差は
\[
 f*_{\pi'_h}g-f*_{\pi_h}g
 =\frac{h^m}{2}D(df,dg)\pmod{h^{m+1}}.
\]
同じ二次関係を両標的で評価すると、交換子の差は
$h^mD(dz_i,dz_j)$ である。
関係式の右辺には既に $h$ が掛かっているため、その差は $h^{m+1}$ に入る。
よって全ての $i,j$ について $D(dz_i,dz_j)=0$ であり、$D=0$ となる。
これは任意の最初の相違を排除する。
$B_h$ は次数一で生成されるので、tensor が等しければ $\phi$ も等しい。
座標変換した模型も同じ生成元固定条件を満たすため、最後の共変性は一意性から従う。
\end{proof}

この議論は構成的である。$\pi_{<m}=\sum_{j<m}h^j\pi_j$ が得られたとき、
\begin{equation}
\label{eq:nfm-normalized-tensor-recursion}
 \pi_m(dz_i,dz_j)=
 -[h^m]\left(
 z_i*_{\pi_{<m}}z_j-z_j*_{\pi_{<m}}z_i
 -h\operatorname{sym}_{*_{\pi_{<m}}}(q_{ij})\right)
\end{equation}
で次の係数が定まる。右辺の星積は、この段階ではグラフによる双微分式として評価する。
最終的な Maurer--Cartan 恒等式は、formality による存在と一意性から確認する。
各固定次数で使うグラフは有限個であり、右辺は二次多項式である。
従って $\pi_m$ は、$q_{ij}$ の有限個の係数と固定したグラフ重みの有限式である。

\begin{proposition}[普遍係数上での正則性と基底変更]
\label{prop:nfm-normalized-model-basechange}
$T$ を被約な有限型複素スキームとし、その上に上記の次数別 PBW 二次代数族と
形式中心三次 $\Omega_h$ があるとする。
各複素閉点で固定入力の星積模型が存在すると仮定する。
このとき正規化された $\pi_h,\phi$ および
\[
 W_h=C_{\pi_h}^{-1}\phi(\Omega_h)
\]
は、各 $h$ 係数が $T$ 上正則な形式族として一意に構成される。
その形成は $T$ からの任意の基底変更と両立する。
基底変更先は非被約でも非平坦でもよい。
\end{proposition}

\begin{proof}
局所枠で\eqref{eq:nfm-normalized-tensor-recursion}を使う。
各係数は有限式なので正則であり、前補題の共変性により局所枠の間で貼り合わさる。
各閉点で再帰は formality の正規化された解と一致する。
従って $[\pi_h,\pi_h]=0$ の各係数は、被約有限型複素スキームの全閉点で零となり、
$T$ 上で零である。これで得られた tensor の星積は結合的である。
再帰により二次関係が成立するので、生成元を固定する代数準同型 $\phi$ が存在する。
各内部次数でこれは剰余が恒等な有限自由加群間の行列であり、
$h$ 進幾何級数によって逆を持つ。

$C_{\pi_h}$ は次数三の有限局所自由加群上で $1+O(h)$ である。
従ってその逆も係数wise正則であり、$W_h$ が正則に得られる。
中心性と formality の接写像の chain 恒等式から
$[\pi_h,W_h]=0$ となる。これは固定入力の補題\texttt{cft:lem:center}と同じ
$h$ 進の HKR 単射性の議論である。
以上は全て有限係数の恒等式と $h$ 進の行列逆であるため、
任意の基底変更後も成立する。
ここで非被約基底について用いたのは、既に構成した恒等式の基底変更であり、
その閉点だけを調べて恒等式を証明したのではない。
\end{proof}

\subsubsection{中心一次空間と一ブロック support の保存}

$U_R=\langle u_1,\ldots,u_r\rangle\subset V_R$ を中心一次空間とし、
$V_R=\langle x,y,z\rangle\oplus U_R$、$P=\partial_x\wedge\partial_y\wedge\partial_z$ とする。
この分解は証明用の局所座標であり、結論はその選択に依存しない。

\begin{lemma}[線型中心の厳密な保存]
\label{lem:nfm-formal-linear-center}
上の生成元固定模型で $u\in U_R$ が $B_h$ の中心に属するなら、
$\iota_{du}\pi_h=0$ である。
さらに任意の $f\in\operatorname{Sym}U_R[[h]]$ について
\[
 C_{\pi_h}(f)=f,
 \qquad f*_{\pi_h}a=fa=a*_{\pi_h}f
\]
が成立する。従って補正後の tensor も、元の中心変数方向を微分しない。
\end{lemma}

\begin{proof}
$h$ の次数について帰納法を使う。
$h^m$ 未満の全 tensor 係数が $U_R$ を消すと仮定する。
これらの係数のみを使う正の次数のグラフで、境界入力に線型関数 $u$ を置く。
$u$ へ辺が入れば、その辺の添字は $U_R$ 方向であり、
出発頂点の tensor のその成分は零である。
辺が一本も入らない場合のグラフの総寄与は、単位の正規化によって零である。
従って $u* a-a*u$ の $h^m$ 係数に寄与するのは
HKR 項 $\pi_m(du,da)$ だけである。
$\phi$ は $u$ を固定し、$u$ は中心元なのでこの交換子は零である。
$a=z_i$ として全成分の消滅を得る。帰納法により $\pi_h$ 全体が相対的となる。

これ以降、$f(U)$ への全 active 微分は零である。
星積の正次数グラフも同じ理由で零となる。
中心写像のグラフで関数頂点 $f(U)$ に辺が入れば零であり、
一本も入らなければその関数頂点は孤立する。
孤立頂点を忘れると積分形式は必要な次元より低い配置空間からの引戻しとなるため、
正次数の重みは零である。これは $C_{\pi_h}(1)=1$ のグラフによる証明でもある。
よって $C_{\pi_h}(f)=f$ を得る。
\end{proof}

\begin{lemma}[三次 Casimir による相対二次 tensor の決定]
\label{lem:nfm-relative-casimir-tensor}
$W_h\in R[[h]][x,y,z,u_1,\ldots,u_r]_3$ とし、
その active 偏微分 $W_x,W_y,W_z$ の Koszul 複体が普遍的に完全であるとする。
二次係数の相対 bivector $Q$ が $\iota_{dW_h}Q=0$ を満たすなら、
一意な $a\in R[[h]]$ によって
\[
 Q=a\,\iota_{dW_h}P
\]
と書ける。$a$ は座標変数に依存しない。
\end{lemma}

\begin{proof}
$Q=Q^{yz}\partial_y\wedge\partial_z+
Q^{zx}\partial_z\wedge\partial_x+Q^{xy}\partial_x\wedge\partial_y$ とし、
$\alpha=Q^{yz}dx+Q^{zx}dy+Q^{xy}dz$ と置く。
Casimir 条件は $d_KW_h\wedge\alpha=0$ である。
三つの偏微分による wedge Koszul 複体の一次完全性から
$\alpha=a\,d_KW_h$ を得る。
$\alpha$ と $d_KW_h$ の係数次数はいずれも二なので、
次数零成分を取れば $a$ は $R[[h]]$ に属する。
次数零の Koszul 単射性により一意である。
全ての座標を内部次数一と数えているため、$a$ に中心変数の依存も残らない。
\end{proof}

active 制限 $g=F|_{U=0}$ が滑らかな平面三次を定める場合、
必要な Koszul 完全性は、active 勾配による八次有限局所自由代数とその次数別分裂完全性から従う。
$h$ 進変形については各固定内部次数で持上げ、完備性により極限を取る。
この議論は、その普遍係数模型からの任意の基底変更と両立する。

\begin{theorem}[残余三次に適合する形式 Poisson 対]
\label{thm:nfm-residual-formal-poisson}
良い一ブロックの PBW 中心三次族に対し、上の固定 formality による
生成元固定模型を取る。量子関係の先頭を
$[a,b]=h\operatorname{sym}\Pi_F(a,b)$、$\Pi_F=\iota_{dF}P$ とし、
$\Omega_h\bmod h=F$ とする。
$V(F)$ と active な平面三次がともに基底上滑らかなら、一意な形式スカラー $a_h=h+O(h^2)$ が存在して
\begin{equation}
\label{eq:nfm-residual-formal-factor}
 \pi_h=a_h\,\iota_{dW_h}P,
 \qquad W_h=C_{\pi_h}^{-1}\phi(\Omega_h)
\end{equation}
となる。従って $\pi_h/h$ の線は $W_h$ と元の rank 三 support によって決まり、
$h$ を逆にした後の $(V(W_h),\pi_h|_{V(W_h)})$ は非零の一ブロック Poisson 対である。

固定入力の完備 MF 移送と相対 window 比較が適用できる場合には、同じ $W_h$ に対し
\[
 \mathcal K_{B_h/(\Omega_h)}\simeq\operatorname{Ku}(V(W_h))
\]
を得る。従ってこの形式 Poisson 対では、残余比較の対象となる三次超曲面が一致する。
この圏同値は入力の結論を適用したものであり、本節で完備 MF 定理の全証明を
独立に再検証したという意味ではない。
\end{theorem}

\begin{proof}
中心逆像の構成から $[\pi_h,W_h]=0$ であり、前補題から $\pi_h$ は相対 bivector である。
Koszul 補題を適用して\eqref{eq:nfm-residual-formal-factor}を得る。
$\pi_h\equiv h\Pi_F\pmod{h^2}$ と $W_h\bmod h=F$ を比較すると
$a_h\equiv h\pmod{h^2}$ である。
従って $a_h/h$ は単元であり、$h$ の逆元を許すと tensor は非零となる。
Jacobi 恒等式は formality 模型の性質から従うが、
相対三変数の Jacobian tensor の表示からも直接従う。
$W_h$ は滑らかな $F$ の形式変形なので、対応する各 Artin 切断とその形式的極限は滑らかな射影三次の変形である。
最後は $W_h$ の定義が固定入力の完備 MF 移送で用いる中心逆像そのものであることによる。
\end{proof}

$\pi_h/h$ の次数一 Poisson 中心は厳密に $U_R[[h]]$ である。
実際、余分な active 線型中心があれば、その定数 active 勾配が
$d_KW_h$ と wedge 零となる。
同じ Koszul 完全性は、それが係数次数 $-2$ の倍数であることを要求するため零である。
従って元の中心一次空間と rank 三 support は、形式変形の途中でも保存される。

\begin{corollary}[形式中心と Casimir 空間の線型同定]
\label{cor:nfm-formal-center-casimir}
上の仮定の下で、補正された $\pi_h$ は $B_h$ だけから定まり、
選んだ中心三次には依存しない。中心逆像は線型同型
\[
 C_{\pi_h}^{-1}\phi:
 Z(B_h)_3\xrightarrow{\sim}
 \operatorname{Cas}(\pi_h)_3
 =R[[h]]W_h\oplus\operatorname{Sym}^3 U_R[[h]]
\]
を与え、純中心三次の部分加群上で恒等である。
特に $\Omega_h$ を $b\Omega_h+G(U)$ に替えると、
$W_h$ は $bW_h+G(U)$ に変わり、$\pi_h$ は変わらない。
ここでは $h$ 進形式中心を扱っており、$h=0$ での通常の中心や、
有限 Artin 切断で新たに現れる $h$-torsion の中心元と同一視しない。
\end{corollary}

\begin{proof}
$\pi_h,\phi$ の構成は二次関係だけを使う。
formality の接写像の chain 恒等式と、その vector 場上での最低次の単射性により、
$C_{\pi_h}$ は Casimir と星積中心の間の双射である。
次数を保つので次数三に制限でき、純中心三次上の恒等性は
補題\ref{lem:nfm-formal-linear-center}による。
$h$ は $R[[h]]$ 上非零因子、$a_h/h$ は単元なので、
三次 $H$ が Casimir である条件は
$d_KH\wedge d_KW_h=0$ である。
Koszul 完全性と次数により $d_KH=c\,d_KW_h$、$c\in R[[h]]$ を得る。
標数零であるため $H-cW_h$ は active 変数に依存せず、純中心三次になる。
和が直和であることも active 勾配の Koszul 単射性から従う。
\end{proof}

\subsubsection{三次 tensor のスカラー消滅}

命題\ref{prop:nfm-relative-residual-second-jet}と同じ中心元の正規化を使い、
$F$ の係数を $h$ に関して定数とする。
$\mathscr H_F=\det\operatorname{Hess}_{x,y,z}(F)$ と書く。
二次の星積の対称性と既知の残余二次 jet を
\eqref{eq:nfm-residual-formal-factor}へ代入すると、
$a_h=h(1+c_2h^2+O(h^3))$ かつ
\begin{equation}
\label{eq:nfm-corrected-poisson-third}
 \pi_h=h\Pi_F+h^3\left(\frac1{24}\Pi_{\mathscr H_F}
                  +c_2\Pi_F\right)+O(h^4)
\end{equation}
となる。Casimir 条件に加えて普遍係数の多項式性と不変式を使うことで、
この $c_2$ を決定できる。

\begin{lemma}[三変数三次形式の二次不変式]
\label{lem:nfm-no-quadratic-cubic-invariant}
三変数三次形式の係数に関する斉次二次の
$\operatorname{SL}_3$ 不変多項式は零である。
\end{lemma}

\begin{proof}
$g=\sum_{i+j+k=3}a_{ijk}x^iy^jz^k$ と書く。
行列式一の対角行列に対する重みを比較すると、不変な斉次二次式は
\[
 Q=Aa_{210}a_{012}+Ba_{201}a_{021}
   +Ca_{120}a_{102}+Da_{111}^2
\]
という形に限られる。実際、二つの係数の添字の和は
$(2,2,2)$ でなければならない。
せん断 $x\mapsto x+ty$ と $x\mapsto x+tz$ に対する
$t=0$ での微分を、それぞれ $L_y,L_z$ とする。
\[
\begin{aligned}
 L_yQ={}&3Aa_{300}a_{012}
       +(A+2C)a_{210}a_{102}+(B+4D)a_{201}a_{111},\\
 L_zQ={}&3Ba_{300}a_{021}
       +(B+2C)a_{201}a_{120}+(A+4D)a_{210}a_{111}.
\end{aligned}
\]
これらはともに零である。まず $A=B=0$ を得て、続いて
$C=D=0$ を得る。
\end{proof}

\begin{proposition}[三次の補正スカラーの消滅]
\label{prop:nfm-third-tensor-scalar-zero}
命題\ref{prop:nfm-relative-residual-second-jet}の正規化と
定理\ref{thm:nfm-residual-formal-poisson}の仮定を用いる。
$F$ の係数は $h$ に依存せず、二次関係を厳密に
$[z_i,z_j]=h\operatorname{sym}(\Pi_F(z_i,z_j))$ とする。
harmonic-angle formality による生成元固定模型では
\begin{equation}
\label{eq:nfm-third-tensor-no-scalar}
 \pi_h=h\Pi_F+\frac{h^3}{24}\Pi_{\mathscr H_F}+O(h^4),
 \qquad a_h=h+O(h^4).
\end{equation}
特に\eqref{eq:nfm-corrected-poisson-third}の $c_2$ は零である。
この恒等式は普遍係数模型からの任意の基底変更と両立する。
\end{proposition}

\begin{proof}
まず\eqref{eq:nfm-normalized-tensor-recursion}を普遍係数で評価する。
関係の係数が $h$ に依存しない場合、$\pi_m$ はその係数に関する
斉次次数 $m$ の多項式である。
これは帰納法による。星積の $h^m$ 係数の各グラフでは、
tensor 入力の $h$ 次数の和が $m$ である。
右辺 $h\operatorname{sym}_*(\Pi_F(z_i,z_j))$ の寄与では、
tensor 入力の次数の和が $m-1$ であり、さらに関係の係数が一つ入る。
従ってどちらの寄与も係数次数 $m$ を持つ。
固定次数でグラフは有限個であり、再帰に係数による除算はない。
またこの多項式構成は線型座標変更に共変である。

中心変数を含まない普遍三次形式
\[
 g=\sum_{|\alpha|=3}a_\alpha x^\alpha,
 \qquad A=\C[a_\alpha:|\alpha|=3],
 \qquad P=\partial_x\wedge\partial_y\wedge\partial_z
\]
を考え、
\[
 D=\pi_3(g)-\frac1{24}\Pi_{\det\operatorname{Hess}(g)}
\]
と置く。$D$ の係数は $A$ の斉次三次式である。
滑らかな三次形式の稠密開では、既存の Casimir の分解と二次 jet により
$D=c\Pi_g$ である。この主張は、必要なら
$g+\sum_j u_j^3$ という滑らかな分離型模型を用いて適用できる。

$v$ を $\Pi_g$ の二次単項式係数を並べたベクトル、$d$ を
$D$ の同じ係数ベクトルとする。
$d_i v_j-d_j v_i$ は稠密開で零なので $A$ 内で恒等的に零である。
従って $\operatorname{Frac}(A)$ 上で $d=cv$ となる。
ここで $c=p/q$ を互いに素な分子・分母で書くと、$cv_i\in A$ より
$q$ は全ての $v_i$ を割る。
一方、$v$ の非零成分は普遍係数 $a_\alpha$ の非零整数倍であり、
全ての $a_\alpha$ が現れる。従ってそれらの最大公約元は $1$ である。
$A$ は一意分解整域なので $q$ は単元、すなわち $c\in A$ を得る。
この分母の除去は、判別式の補集合上の正則性だけからは従わない点に注意する。

斉次性から $c$ は係数次数二である。
$P$ は $\operatorname{SL}_3$ に不変であり、再帰と Hessian は共変なので、
$d=cv$ におけるスカラーの一意性から $c$ は
$\operatorname{SL}_3$ 不変である。
補題\ref{lem:nfm-no-quadratic-cubic-invariant}により $c=0$ となる。

次に中心変数と任意の混合係数を許す。
補題\ref{lem:nfm-formal-linear-center}により全ての tensor 入力は
active 方向だけを持つ。従ってグラフと再帰の全微分は active 方向であり、
$U=0$ への制限と可換である。
$g=F|_{U=0}$ とすれば
\[
 \pi_3(F)|_{U=0}=\pi_3(g),\qquad
 \mathscr H_F|_{U=0}=\det\operatorname{Hess}(g).
\]
既存の分解
$\pi_3(F)-\Pi_{\mathscr H_F}/24=c_2(F)\Pi_F$ のスカラーは
中心変数を含む全座標に依存しない。
制限を取ると $c_2(F)\Pi_g=0$ となり、active 勾配の
次数零 Koszul 単射性から $c_2(F)=0$ を得る。
以上の普遍恒等式を基底変更すれば、非被約基底についても同じ結論を得る。
最後に $W_h=F+h^2\mathscr H_F/24+O(h^3)$ と
$\pi_h=a_h\iota_{dW_h}P$ を比較すれば $a_h=h+O(h^4)$ となる。
\end{proof}

\begin{remark}[規格化と適用範囲]
この証明が用いる中心元の情報は二次 jet だけである。
高次の中心元を Hesse 座標で記述した際の選択は、
$D=\pi_3-\Pi_{\mathscr H}/24$ の普遍多項式性と共変性に影響しない。
一方、一般の formality への変更、$h$ の再パラメータ化、
または $h$ に依存する二次関係を追加した場合には、
同じ係数式をそのまま主張しない。

偶奇性を用いた剰余の改善は、以下の系\ref{cor:nfm-mixed-even-jet}で証明する。
\end{remark}

\subsubsection{反対積による全次数の偶奇性}

ここでは harmonic-angle の Kontsevich formality を固定する。
一般の angle form や、任意に gauge 変換した formality に同じ偶奇性を
無条件に主張するものではない。
\begin{lemma}[harmonic-angle グラフの反転]
\label{lem:nfm-harmonic-parity}
同じ bivector $\pi$ を $k$ 個入れた星積の項を $B_k(\pi;f,g)$ と書く。
このとき
\[
 B_k(\pi;g,f)=(-1)^kB_k(\pi;f,g),\qquad
 f*_{-\pi}g=g*_{\pi}f.
\]
さらに関数上の接 formality 写像について
\[
 C_{-\pi}(f)=C_{\pi}(f)
\]
が成立する。最後の等式は各奇数 $k$ に対して
$\mathcal U_{k+1}(f,\pi,\ldots,\pi)=0$ であることから従う。
\end{lemma}

\begin{proof}
Kontsevich の配置空間と向きの規約を用いる
\cite[Sections~2.2 and~6.2]{NFMKontsevich}。
上半平面の反射 $r(z)=-\overline z$ は harmonic-angle 一形式を
$r^*d\varphi=-d\varphi$ に変える。
星積では境界の二点を $0,1$ に固定した表示を使い、
$z\mapsto1-\overline z$ により両境界点のラベルを交換する。
内部点は $k$ 個なので配置空間の向きは $(-1)^k$ 倍される。
辺は $2k$ 本であり角形式の積には符号が生じない。
従って両入力を交換したグラフの重みは $(-1)^k$ 倍になり、第一式を得る。
多重線型性から第二式も従う。

中心写像の $k$ 次項では、内部点は bivector の $k$ 頂点と関数の一頂点の
合計 $k+1$ 個、境界点は零個、辺は $2k$ 本である。
配置空間 $C_{k+1,0}$ の次元は $2k$ である。
商を取る前の配置空間の向きは反射により $(-1)^{k+1}$ 倍される一方、
実平行移動と正の拡大の二次元群の向きは $-1$ 倍される。
従って商の向きは $(-1)^k$ 倍される。
ここでは境界ラベルの交換はなく、同じグラフに戻る。
角形式は $2k$ 個なのでその積は不変であり、
重み $w_\Gamma$ は $w_\Gamma=(-1)^k w_\Gamma$ を満たす。
標数零では奇数 $k$ の重みは零である。
これは二次 tensor の線型性・二次性を仮定しないグラフの対称性であり、
偶数次数の中心グラフ全てが消えるという主張ではない。
\end{proof}

\begin{proposition}[生成元固定模型の偶奇性]
\label{prop:nfm-normalized-parity}
補題\ref{lem:nfm-normalized-star-model}の存在仮定の下で、
二次関係の $q_{ij}$ は $h$ に依存しないとする。
すると正規化した tensor は
\[
 \pi_{-h}=-\pi_h
\]
を満たす。Weyl 対称化を $\sigma_h$ とすれば、線型写像
$T_h=\phi_h\sigma_h$ も $T_{-h}=T_h$ を満たす。
特に $\Omega_h=\sigma_h(F_h)$、$F_{-h}=F_h$ という規格化の形式中心元に対し、
\[
 W_h=C_{\pi_h}^{-1}\phi_h(\Omega_h)
 \quad\text{は}\quad W_{-h}=W_h
\]
を満たす。
\end{proposition}

\begin{proof}
生成元を固定して語順を逆にする写像は
$B_{-h}\simeq B_h^{\mathrm{op}}$ を与える。
前補題により $*_{-\pi_h}$ は $*_{\pi_h}$ の反対積である。
従って $B_{-h}$ の生成元固定模型を二通りに書き、
補題\ref{lem:nfm-normalized-star-model}の一意性を適用すれば
$\pi_{-h}=-\pi_h$ となる。

Weyl 対称化された語の総和は語順を逆にしても不変である。
このため同じ反対積の比較で $T_{-h}=T_h$ を得る。
また $C_{\pi_{-h}}=C_{-\pi_h}=C_{\pi_h}$ なので、
$F_h$ が偶数級数なら $C_{\pi_h}^{-1}T_h(F_h)$ も偶数級数になる。
全ての等式は各係数の有限グラフ恒等式であり、
命題\ref{prop:nfm-normalized-model-basechange}で得た普遍族からの
任意の基底変更後も成立する。
\end{proof}

\begin{corollary}[全混合係数における jet の改善]
\label{cor:nfm-mixed-even-jet}
命題\ref{prop:nfm-relative-residual-second-jet}の全混合 Hesse 族で、
同じ規格化中心元 $\widetilde\Omega_{F,h}$ を用いる。
このとき
\[
 W_h=F+\frac{h^2}{24}\mathscr H_F+h^4W_4(F)+O(h^6),
 \qquad
 W_h\equiv H^\sharp\pmod{h^4}.
\]
また形式スカラー $a_h$ も奇数級数である。
命題\ref{prop:nfm-third-tensor-scalar-zero}の $c_2=0$ と合わせると
\[
 \pi_h=h\Pi_F+\frac{h^3}{24}\Pi_{\mathscr H_F}+O(h^5)
       \equiv\Pi^\sharp\pmod{h^5}.
\]
これは全 $n\ge3$、全混合係数についての合同である。
次に決定すべき potential の係数は $W_4$ であり、
$h^3$ の混合補正は存在しない。
\end{corollary}

\begin{proof}
指定中心元は
$\sigma_h(F-h^2\mathscr H_F/12+h^4\Psi_\lambda(F))$
と書かれ、括弧内は $h$ の偶数多項式である。
前命題と既知の二次 jet を合わせる。
代数的補正は厳密に
$H^\sharp=F+h^2\mathscr H_F/24+h^4\tau_F(\Psi_\lambda(F))$
であるため、第一の合同が従う。
$\pi_h=a_h\Pi_{W_h}$ に $h\mapsto-h$ を施すと
$(a_{-h}+a_h)\Pi_{W_h}=0$ となる。
相対勾配 Koszul 複体の完全性からスカラーの作用は単射であり、
$a_{-h}=-a_h$ を得る。
最後は $c_2=0$ と tensor の奇数性による。
ここでは $W_4=\tau_F(\Psi_\lambda(F))$ とも、
有限式 $H^\sharp$ が全次数の中心逆像であるとも結論していない。
\end{proof}

\begin{remark}[中心次数の filtration]
\label{rem:nfm-central-filtration}
相対グラフの全ての微分方向は $x,y,z$ であり、中心変数を微分しない。
従って $I=(U)$ による filtration を保ち、$U=0$ への制限は
生成元固定再帰、Weyl 対称化の像、中心写像およびその逆と可換である。
さらに各グラフの中心次数は、各頂点に選んだ係数の中心次数の和である。
従って tensor の係数を $I^d/I^{d+1}$ で求めるために、
元の関係の中心次数 $d$ を超える係数は不要である。
これは $\pi_h$ の境界部分を純三変数の計算から求める根拠となるが、
その情報のみから正の中心次数を持つ全混合係数は決まらない。
中心三次自身の選択に正の中心次数の係数が含まれる場合には、
もちろんその指定係数も入力に含める。
\end{remark}

\subsubsection{五次の尺度補正と代数的処方の適用範囲}

\begin{lemma}[log-canonical 模型の厳密な逆 tensor]
\label{lem:nfm-log-canonical-inverse}
$g=xyz$ とし、二次関係を
$[x,y]=h(xy+yx)/2$ およびその巡回置換とする。
harmonic-angle formality の生成元固定模型では
\[
 \pi_h(g)=\log\left(\frac{1+h/2}{1-h/2}\right)\Pi_g
       =\left(h+\frac{h^3}{12}+\frac{h^5}{80}+O(h^7)\right)\Pi_g.
\]
\end{lemma}

\begin{proof}
Kontsevich の元の星積による
$\sigma=\sum_{i<j}\lambda_{ij}x_ix_j\partial_i\wedge\partial_j$
の量子化では、原座標に関する交換比が
$q_{ij}(t\sigma)=\exp(t\lambda_{ij})$ となる
\cite[Section~1.6; Theorem~6.9 and Remark~6.10, pp.~59--60]{NFMLindbergPym}。
ここで交換比は、トーラス同変性による
$x_i\star x_j=w_{ij}x_ix_j$ の係数の比 $w_{ij}/w_{ji}$ である。
従って抽象的な量子トーラスの同型による生成元変更は不要である。

$q_h=(1+h/2)/(1-h/2)$、$\ell(h)=\log q_h$ と置き、
この恒等式を $t\mapsto\ell(h)$ で形式的に基底変更する。
$\ell(h)\Pi_{xyz}$ の星積は、原生成元で
$x\star y=q_h y\star x$ とその巡回置換を満たす。
これは指定した $h$ を保った二次関係と同値である。
単元係数の $q$-交換関係には順序単項式の PBW 基底があり、
星積側も同じ単項式の単元倍を基底に持つ。
従って生成元を固定する準同型は、各内部次数で剰余が恒等な
$\C[[h]]$-線型同型である。
補題\ref{lem:nfm-normalized-star-model}の一意性が所望の式を与える。
ここで $\ell(h)$ は標的 tensor の係数であり、
比較の際に指定されたパラメータ $h$ 自体を取り替えたのではない。
\end{proof}

\begin{proposition}[純 active 族の五次 tensor]
\label{prop:nfm-active-fifth-covariant}
普遍三変数三次形式 $g$ に対する生成元固定 tensor を
$\pi_h(g)=\sum_{m\ge1}h^m\pi_m(g)$ とする。
$x_1=x,x_2=y,x_3=z$、$\varepsilon^{123}=1$ と規約を取り、
$t_{ijk}=\partial_i\partial_j\partial_k g/6$ とし、
\[
 S(g)=54\,t_{abc}t_{def}t_{ghi}t_{jkl}
             \varepsilon^{adg}\varepsilon^{bej}
             \varepsilon^{chk}\varepsilon^{fil}
\]
と定める。添字は $1,2,3$ を走り、反復添字について和を取る。
この $S$ は古典的な四次 Aronhold 不変式であり、
四つの volume の縮約による表示も既知である
\cite[Section~3, after Remark~3.3]{NFMOttavianiAronhold}。
ここでは $S(xyz)=1$ と規格化する。
harmonic-angle formality の生成元固定模型では
\[
 \pi_5(g)=\frac1{80}S(g)\Pi_g
\]
となる。
\end{proposition}

\begin{proof}
$\pi_m(g)$ と volume によって双対な一次形式を $\alpha_m(g)$ とする。
滑らかな $g$ に純中心 Fermat 三次を加えた分離型を用いると、
Casimir の分解と $U=0$ への制限により
$\sum_m h^m\alpha_m=a_h\,dW_h$ であり、$a_h$ は座標に依存しない。
従って $d\alpha_m=0$ である。
これは普遍係数に関する多項式恒等式なので、全ての $g$ で成立する。
$E=x\partial_x+y\partial_y+z\partial_z$ に対し
\[
 G_m(g)=\frac13\iota_E\alpha_m(g)
\]
と置くと、Cartan の公式と内部次数により $dG_m=\alpha_m$ となる。
従って $G_m$ は係数次数 $m$、変数次数三の
$\operatorname{SL}_3$ 共変式である。

この型の共変式の次元は
$\operatorname{Sym}^m(\operatorname{Sym}^3\C^3)$ における
最高ウェイト $(m+2,m-1,m-1)$ の重複度に等しい。
重みの個数を
\[
 p_m(a,b,c)=[t^mX^aY^bZ^c]
     \prod_{i+j+k=3}(1-tX^iY^jZ^k)^{-1}
\]
と置く。$\delta=(2,1,0)$ とすれば、Weyl の交代和より
最高ウェイト $\lambda$ の重複度は
\[
 \sum_{\sigma\in S_3}\operatorname{sgn}(\sigma)
             p_m(\lambda+\delta-\sigma\delta)
\]
である。$m=5$、$\lambda=(7,4,4)$ では、この六項は
\[
 41-35-30+23+21-19=1
\]
となる。
一方、上に定義した $S$ は明示的な volume 縮約により不変であり、
$S(xyz)=1$ なので非零である。
従って $g\mapsto S(g)g$ は当該一次元空間を張り、
$G_5(g)=\kappa S(g)g$ を得る。
補題\ref{lem:nfm-log-canonical-inverse}と $S(xyz)=1$ により
$\kappa=1/80$ と決まる。
\end{proof}

\begin{corollary}[混合係数の制限と代数的 tensor 補正の限界]
\label{cor:nfm-fifth-restriction-obstruction}
任意の混合三次 $F$ に対し $g=F|_{U=0}$ と置くと、
\[
 \pi_5(F)|_{U=0}=\frac1{80}S(g)\Pi_g
\]
となる。従って先の Hessian 規格化の族に対する代数的補正
\[
 \pi_h^\sharp=h\Pi_F+\frac{h^3}{24}\Pi_{\mathscr H_F}
\]
は、$S(g)\ne0$ なる稠密開では、固定した $h$ と生成元の下で
実際の形式 tensor $\pi_h$ と $h^5$ の係数から異なる。
この開は全ての $n\ge3$ で非空である。
\end{corollary}

\begin{proof}
線型中心の保存により全グラフの微分は active 方向であるため、
$U=0$ への制限と再帰は可換である。
従って純 active の公式を制限に適用できる。
代数的 tensor 補正の $h^5$ 係数は零であるが、
上の制限は $S(g)\ne0$ で零でない。
五次より低い次数の一致は系\ref{cor:nfm-mixed-even-jet}による。
実際、$g=(x^3+y^3+z^3)/3+xyz$ は滑らかで $S(g)=-7$ であり、
$F=g+\sum_j u_j^3$ は全次元で滑らかな分離型を与える。
普遍係数空間の既約性から非消滅条件は稠密開を定める。
\end{proof}

\begin{remark}[規格化の検算と混合係数]
上の四つの $\varepsilon$ の縮約は $6^4$ 個の有理数項で厳密に展開できる。
得られる $S$ は通常の単項式係数に関する整数係数二十五項であり、
\[
 S\left(\frac{x^3+y^3+z^3}{3}+\lambda xyz\right)
       =\lambda(\lambda^3-8)
\]
を満たす。
共変式の同じ重複度を $m=1,\ldots,8$ で計算すると
$1,0,1,0,1,0,2,0$ となる。
これらは上の有限生成関数による整数計数であり、
Kontsevich グラフ重みや残余同値を計算したものではない。

この結果は純 active 係数についての還元である。
中心変数を含む一般の混合三次について、
$\pi_5(F)=S(F|_{U=0})\Pi_F/80$ と結論することはできない。
active 方向の局所 jet による別の共変式が混合項に寄与する可能性は残る。
また系\ref{cor:nfm-fifth-restriction-obstruction}は固定された tensor の不一致であり、
tensor のスカラーを忘れた比較、残余三次の四次 jet、
または Kuznetsov 成分の非同値までは結論しない。
\end{remark}


\subsubsection{五次の混合項と残余四次項}

次に、五次 tensor の全混合係数を二つの普遍定数へ還元する。
まず共変式の空間を決定し、続いて既知の graph weight を使って二定数を評価する。
active 微分について
\[
 T_{ijk}=F_{ijk},\qquad M_{ij}=F_{ij},\qquad q_i=F_i
\]
と置く。$T$ は座標 $x,y,z,U$ に依存せず、$F|_{U=0}=g$ だけから決まる。
全添字は $1,2,3$ を走り、$\varepsilon^{123}=1$ とする。
次の完全縮約を定義する。
\begin{align*}
 S(T)&=\frac1{24}
 T_{adg}T_{bej}T_{chk}T_{fil}
 \varepsilon^{abc}\varepsilon^{def}\varepsilon^{ghi}\varepsilon^{jkl},\\
 I(F)&=\frac1{24}
 T_{adg}T_{bej}M_{cf}M_{hk}M_{il}
 \varepsilon^{abc}\varepsilon^{def}\varepsilon^{ghi}\varepsilon^{jkl},\\
 J(F)&=\frac1{24}
 T_{adg}T_{bej}T_{chk}M_{fi}q_l
 \varepsilon^{abc}\varepsilon^{def}\varepsilon^{ghi}\varepsilon^{jkl}.
\end{align*}
$S$ は $S(xyz)=1$ の規格化の四次 Aronhold 不変式であり、
$S(g_\lambda)=\lambda(\lambda^3-8)$ である。
$I(F),J(F)$ は全変数に関する三次式である。

\begin{proposition}[五次の混合項の二定数への還元]
\label{prop:nfm-mixed-fifth-reduction}
固定 harmonic-angle formality に依存する二つの普遍定数
$\beta,\gamma\in\mathbf C$ があり、全 $n\ge3$ の全混合係数について
\[
 \pi_5(F)=\Pi_{Q_5(F)},\qquad
 Q_5(F)=\frac1{80}S(g)F
       +\beta\bigl(I(F)-S(g)F\bigr)
       +\gamma\bigl(J(F)-S(g)F\bigr).
\]
ここで $\pi_h=h\Pi_F+h^3\Pi_{\mathscr H_F}/24+h^5\pi_5+O(h^7)$ である。
係数 $1/80$ は別途得た純三変数の五次 tensor の計算を用いる。
この入力を使わない場合も、右辺は三つの普遍定数
$\alpha S(g)F+\beta I(F)+\gamma J(F)$ の勾配へ還元できる。

\end{proposition}

\begin{proof}
生成元固定再帰の $\pi_5$ は、関係 $\Pi_F$ の係数に関する
斉次次数五の微分多項式であり、active 方向だけを微分する。
再帰の Weyl 対称化も局所 jet で書けることを確認しておく。
任意の二次式 $r$ に対して
\[
 \operatorname{sym}_*(r)=r+\frac12\sum_{a,b}
      (\partial_a\partial_b r)(z_a*z_b-z_az_b)
\]
である。$r=\Pi_F(z_i,z_j)$ なら、その active 二階微分は $T$ の成分である。
片方の添字が中心なら括弧内は零となる。
active 添字の場合、正次数の単位正規化グラフは各線型入力をちょうど一回微分するので、
括弧内も下位 tensor の jet だけで書け、裸の座標は残らない。
従って帰納的に再帰の全項が以下の局所 jet 型に入る。
各係数をある点で評価すると、使える jet は $q,M,T$ だけである。
$F$ 自身は二次関係には現れない。
$q,M,T$ の全座標次数はそれぞれ $2,1,0$ であり、
$\pi_5$ の係数は全座標に関する次数二の多項式である。
従って現れ得る型は $qT^4$ と $M^2T^3$ の二つに限られる。
固定体積形式により $\pi_5$ を一形式 $\alpha_5$ と同一視する。
$\operatorname{SL}_3$ 共変性と determinant weight を合わせると、
これら二型の共変式空間の次元はそれぞれ $1,2$ である。

この小さい表現論の計算は、次の有限係数計算で直接確認できる。
$H_{d,k}(X,Y,Z)$ を
\[
 H_{d,k}=[t^k]\prod_{a+b+c=d}(1-tX^aY^bZ^c)^{-1}
\]
と置く。$\lambda$ の重複度は
$\rho=(2,1,0)$ として次の式で与えられる。
\[
 \sum_{w\in\mathfrak S_3}\operatorname{sgn}(w)
 [X^{\lambda_1+\rho_1-(w\rho)_1}
  Y^{\lambda_2+\rho_2-(w\rho)_2}
  Z^{\lambda_3+\rho_3-(w\rho)_3}]H.
\]
二型について $\lambda=(5,4,4)$ を使うと
\[
\begin{array}{c|c|c}
H&\text{六つの符号付き係数}&\text{和}\\\hline
H_{1,1}H_{3,4}&67,-58,-55,41,42,-36&1\\
H_{2,2}H_{3,3}&158,-136,-126,94,92,-80&2
\end{array}
\]
となる。

三つの一形式 $S(g)d_KF,d_KI(F),d_KJ(F)$ はこの三次元空間に属し、
一次独立である。具体的には $T$ を $g_1$ の三階微分とし、
ある点で
\[
 q=(1,0,0),\qquad
 M=\begin{pmatrix}1&1&0\\1&0&0\\0&0&0\end{pmatrix}
\]
とすると、三つの一形式を列に持つ行列は
\[
 \begin{pmatrix}
 -7&-2/3&-10/3\\
 0&-2/3&-1\\
 0&5/3&0
 \end{pmatrix},\qquad \det=-35/3\ne0
\]
となる。ここで $q,M,T$ は独立な Taylor 係数として指定できる。
実際、一つの中心変数 $u_1$ を用いて
$F=T_{ijk}x_ix_jx_k/6+u_1M_{ij}x_ix_j/2+u_1^2q_ix_i+cu_1^3$
とし、$(x,u_1)=(0,1)$ で評価すればよい。
したがってこの独立性は全ての $n\ge3$ に有効である。
したがって三つの勾配は共変式空間全体を張る。
このため $\alpha_5=d_K(\alpha SF+\beta I+\gamma J)$ と書ける。
全 graph の有効添字は三つの active 方向に限られるため、
三定数は中心変数の個数に依存しない。

純三変数の場合、$I(g),J(g)$ はいずれも
$S(g)g$ に等しい。これも小さい共変式計算で確認できる。
$H_{3,5}$ に $\lambda=(7,4,4)$ を適用した六項は
$41,-35,-30,23,21,-19$ で、和は一である。
従って次数五・三次出力の共変式は $S(g)g$ の一次元空間である。
$g=g_1$、$(x,y,z)=(1,0,0)$ で
$S=-7$、$g=1/3$、$I=J=-7/3$ を計算すれば比例定数は一となる。
純 active 計算 $\pi_5(g)=S(g)\Pi_g/80$ から
$\alpha+\beta+\gamma=1/80$ が従い、表示された形を得る。
以上の係数計算と縮約の評価は、保存した検証コードでも有理数として確認できる。
局所 jet の多項式恒等式なので、普遍模型からの任意の基底変更とも両立する。
\end{proof}

\begin{remark}[分離型で消える二方向]
\label{rem:nfm-mixed-fifth-potential}
$I-SF$ と $J-SF$ は $U=0$ に制限すると零となる。
より一般に $F=g(x+a(U))+G(U)$ という中心平行移動された分離型では、
両者は純中心三次 $-S(g)G(U)$ に等しく、その active 勾配は零である。
従ってこの分離型の計算だけでは二つの普遍定数を区別できない。
以下では二つの混合 Taylor jet で定数を決定し、
さらに中心逆像の局所性から四次 potential とその幾何的比較類を求める。
\end{remark}

% !TEX root = ../nc-cubic.tex

\begin{proposition}[混合五次 tensor の普遍定数]
\label{prop:nfm-mixed-fifth-explicit}
命題\ref{prop:nfm-mixed-fifth-reduction}の harmonic-angle 規格化において
\[
 \beta=\frac1{60},\qquad \gamma=\frac1{240}.
\]
従って全ての $n\ge3$ の混合三次 $F$ に対し
\[
 \boxed{\quad\pi_5(F)=\frac1{240}\Pi_{\,4I(F)+J(F)-2S(g)F}\quad}
\]
となる。この式は普遍係数の多項式恒等式であり、既存の普遍 PBW 模型からの
任意の基底変更と両立する。
\end{proposition}

\begin{proof}
先の二定数への還元により、二つの独立な局所 jet で正規化再帰を評価すれば足りる。
一つの中心変数 $u$ を用いた斉次三次
\[
 F_A=xyz+\frac12uz^2,\qquad F_B=xyz+u^2z
\]
を採用する。$u=1$ とし、$p=(x,y,z)=(0,0,1)$ で評価する。
いずれも $g=xyz$、$S(g)=1$ であり、完全縮約の定義から
\[
\begin{array}{c|cc}
 &I(F)-S(g)F&J(F)-S(g)F\\\hline
 F_A&-2uz^2/3&-uz^2/6\\
 F_B&-u^2z&-2u^2z/3
\end{array}
\]
を得る。従って求める二つの成分は
\[
 \pi_5(F_A)^{xy}(p)=\frac1{80}-\frac{4\beta+\gamma}{3},\qquad
 \pi_5(F_B)^{xy}(p)=\frac1{80}-\beta-\frac23\gamma.
\]
ここでは特異な二つの三次を独立 jet の評価に用いている。
一般対象の主張を特異な例から直接推定するのではなく、
既に証明した普遍多項式の二定数表示に代入している。

$P=\Pi_F$、$R=\Pi_{\det\operatorname{Hess}_K(F)}/24$ とする。
$B_k(P^{\otimes k};f,g)$ を、先頭項が $fg+\frac12 P(df,dg)$ となる
Kontsevich 星積の $k$ 頂点部分とし、$B_k[P^{k-1},R]$ は
各頂点への $R$ の挿入の和を表す。
$q^{ij}=P(dx_i,dx_j)$ とすると、生成元固定再帰の五次部分は
\begin{align*}
 \pi_5^{ij}
 ={}&-2\left(B_5[P^5;x_i,x_j]+B_3[P^2,R;x_i,x_j]\right)\\
 &+\frac12\sum_{a,b=1}^3 \partial_a\partial_b q^{ij}
 \left(B_4[P^4;x_a,x_b]+B_2[P,R;x_a,x_b]\right).
\end{align*}
偶奇性により奇数頂点部分は交換子に二倍で現れ、
右辺の最後の項は二次式の Weyl 対称化から生じる。

この有限式の評価には、Banks--Panzer--Pym\cite[Corollary~1.3; Section~6]{NFMBanksPanzerPym}の harmonic weights を
Buring\cite[Sections~3.4--3.6]{NFMBuringGraphs}がまとめた \texttt{star5.txt} と、同じ著者資料の
\texttt{kontsevich\_2\_k.txt},
\texttt{weights\_kontsevich\_2\_k.txt} ($1\le k\le5$) を用いる。
規約は Buring の学位論文 Section~3.4.1--3.4.2, pp.~74--76 にある。
前者の $k$ 次係数は $fg+hP(df,dg)$ の規約であり、
本稿の規約では $2^k$ で割る。
頂点置換と outgoing edge の交換をまとめた multiplicity は
$m(\Gamma)=2^k k!/|\operatorname{Aut}\Gamma|$ なので、
元の weight $w(\Gamma)$ に対する本稿の係数は
$w(\Gamma)/|\operatorname{Aut}\Gamma|$ となる。
この二種類のデータの一致も、有理数演算と有限置換の列挙で確認できる。

線型入力では二つの sink の入次数がともに一でなければ項が消え、
二次 tensor では aerial vertex の入次数が三以上なら項が消える。
これらの条件で残る非零項数は $k=1,2,3,4,5$ の順に
$1,1,1,19,51$ である。
上の式の三項の評価は次の表になる。
最後の列 $W$ は $\frac12\sum_{a,b}\partial_{ab}q^{xy}(B_4+B_2)$ を表す。
\[
\begin{array}{c|ccc|c}
 F&B_5^{xy}&B_3[P^2,R]^{xy}&W&-2(B_5+B_3)+W\\\hline
 F_A&1/480&0&-1/144&-1/90\\
 F_B&1/720&0&-1/240&-1/144
\end{array}
\]
全ての項は保存した73項の有限 graph certificate と\texttt{verify\_mixed\_fifth.py}で有理数として評価される。
この計算は既知の graph weight を入力とし、その積分自体を再計算してはいない。
元の graph--weight 対との符号・自己同型群を含む照合、
公刊の二変数 log-canonical 星積の五次までの係数との照合も行う。
従って
\[
 4\beta+\gamma=\frac{17}{240},\qquad
 \beta+\frac23\gamma=\frac7{360},
\]
であり、これを解くと $\beta=1/60$、$\gamma=1/240$ を得る。
\end{proof}


\subsubsection{中心逆像の四次項と一つの幾何的障害}

以下では指定した Hessian 正規化中心元を
\[
 \widetilde\Omega_h=\operatorname{sym}
 \left(F-\frac{h^2}{12}\mathscr H_F+h^4\Psi_\lambda(F)\right)
\]
とする。$S,I,J$ の規約は命題\ref{prop:nfm-mixed-fifth-reduction}と同じである。
この中心元の規格化は四次のスカラーを決めるために必要であり、
任意の単元倍に対して同じスカラー公式を主張するものではない。

\begin{lemma}[四次中心逆像の局所共変式]
\label{lem:nfm-fourth-center-locality}
固定した harmonic-angle 模型に依存する二定数 $d,e\in\C$ があって、
全混合係数に対して
\[
 W_4(F)=\Psi_\lambda(F)+d I(F)+e J(F)
\]
となる。二定数は $n\ge3$ に依存しない。
\end{lemma}

\begin{proof}
$R_{\pi_h}=\phi\circ\operatorname{sym}$ と置く。
これは可換多項式を同じ生成元の星積で完全対称化する写像である。
Kontsevich グラフと完全対称化は active な affine 座標変更に共変であり、
$R_{\pi_h}(1)=C_{\pi_h}(1)=1$ である。
したがって正次数の係数をある点で評価すると、
入力関数を少なくとも一回微分した局所 jet の多項式で表せる。
具体的には、その点を原点へ移して入力多項式を Taylor 展開すればよい。
定数項への作用は単位の保存で零となり、
残る有限個の単項式の星積は有限個のグラフで計算される。
グラフの微分添字は全て active 方向であるので、
純中心関数への正次数の作用も零である。

$h^4$ の係数には、$R$ と $C^{-1}$ の四次項のほか、
$\pi_3=\Pi_{\mathscr H_F}/24$ を二次項へ入れた項と、
$-h^2\mathscr H_F/12$ への二次作用が含まれる。
いずれも $F$ の係数次数は五、全座標次数は三である。
中心元の $h^4\Psi_\lambda(F)$ 自身への正次数の作用は
$h^4$ の係数には寄与しない。
従って $W_4-\Psi_\lambda(F)$ は、
$q_i=F_i$, $M_{ij}=F_{ij}$, $T_{ijk}=F_{ijk}$
だけの局所多項式であり、微分されない $F$ は現れない。
次数から可能な型は $qMT^3$ と $M^3T^2$ に限られる。

この二型の $\operatorname{SL}_3$ スカラー共変式の空間はいずれも一次元である。
命題\ref{prop:nfm-mixed-fifth-reduction}の Weyl 係数公式を
最高重み $(4,4,4)$ に適用すると、六項はそれぞれ
\[
 \begin{aligned}
 H_{1,1}H_{2,1}H_{3,3}&:
       &&156,-136,-136,100,109,-92,\\
 H_{2,3}H_{3,2}&:
       &&122,-105,-105,78,84,-73
 \end{aligned}
\]
で、和はともに一である。
非零の $J,I$ がそれぞれを張るので表示を得る。
有効な微分添字は三つの active 方向に限られ、二定数は $n$ に依存しない。
\end{proof}

\begin{lemma}[Hessian 写像と Aronhold 不変式による表示]
\label{lem:nfm-classical-mixed-covariants}
$\mathscr H(F)=\det\operatorname{Hess}_K(F)$ とし、
$T=D_K^3F$, $M=D_K^2F$, $q=d_KF$ と置く。
$S(T)$ は上の規格化の Aronhold 四次式である。
このとき
\[
 \begin{aligned}
 I(F)
 &=\frac1{12}\operatorname{tr}
       \bigl(\operatorname{adj}(M)\operatorname{Hess}_K(\mathscr H(F))\bigr)
   =\frac1{12}D\mathscr H_F[\mathscr H(F)],\\
 J(F)&=\frac14(dS)_T[M\odot q],\qquad
 (M\odot q)_{ijk}
       =\frac{M_{ij}q_k+M_{ik}q_j+M_{jk}q_i}{3}.
 \end{aligned}
\]
第二式の微分は対称三階 tensor $T$ の空間上で取り、
$M\odot q$ の対称化には表示した平均を用いる。
従って $I,J$ は Hessian 写像の微分と Aronhold 四次式の微分によって
計算でき、独立した新しい不変式を導入する必要はない。
\end{lemma}

\begin{proof}
第一式の trace は $M^3T^2$ 型、第二式の右辺は $qMT^3$ 型の
スカラー共変式である。
それぞれの共変式空間は前の計数により一次元なので、
$I,J$ との比例定数を一つの非零評価で決められる。
$F=xyz$ と $(x,y,z)=(1,1,1)$ を取る。
このとき $\mathscr H(F)=2F$ であり、
\[
 \operatorname{tr}\bigl(\operatorname{adj}(M)
           \operatorname{Hess}(\mathscr H(F))\bigr)
       =2\operatorname{tr}(\operatorname{adj}(M)M)=12.
\]
一方 $I=1$ である。
また $T_{123}=1$ と添字対称性だけが非零で、
$q=(1,1,1)$、$M_{ii}=0$, $M_{ij}=1$ ($i\ne j$) となる。
$S(T)=1$、$(M\odot q)_{123}=1$ であり、
四次式を微分すると $(dS)_T[M\odot q]=4$、$J=1$ となる。
これで両比例定数が定まる。
Hessian 写像の方向微分の表示は、行列式の微分公式による。
\end{proof}

体積を局所的に固定しない表示では、同じ公式で
$\mathscr H$ を $\mathscr H_P$、$S$ を $S_P$ に置き換え、
方向微分では $P$ を固定する。
これは各式の四つの $P$ の重みと一致する。

\begin{lemma}[Hesse 中心元の四次恒等式]
\label{lem:nfm-fourth-center-symbol}
\[
 F=g_\lambda+Axy+Byz+Czx+Dx+Ey+Hz+G
\]
とし、$A,B,C$ は中心一次式、$D,E,H$ は中心二次式、$G$ は中心三次式とする。
$S=\lambda(\lambda^3-8)$ と書く。
指定中心元の四次 Weyl symbol は恒等的に
\[
 \Psi_\lambda(F)=-\frac{S}{80}F-\frac1{40}I(F)
                         +\frac1{60}J(F)+\mathcal C_4(U),
\]
ただし
\[
 \mathcal C_4(U)=
  \frac{\lambda}{60}(A^3+B^3+C^3)
  +\frac{\lambda^2}{40}ABC
  +\frac{2-\lambda^3}{180}(AH+BD+CE)
  +\frac{S}{80}G
\]
である。
\end{lemma}

\begin{proof}
固定入力論文A\cite{NCCubicA}の
系\texttt{cor:pg-hesse-hessian-form}、式\texttt{eq:pg-hesse-hessian-closed}
にある有限中心公式を $h^4$ について展開する。
$Q=Axy+Byz+Czx$, $Z_2=Az^2+Bx^2+Cy^2$,
$L=Dx+Ey+Hz$, $R_2=B^2x+C^2y+A^2z$,
$R_{11}=ACx+ABy+BCz$ と置けば、その係数は
\[
 \Psi_\lambda(F)=
 -\frac S{48}g_\lambda-\frac{\lambda^4}{48}Q
 +\frac{2-\lambda^3}{36}Z_2-\frac S{144}L
 -\frac{\lambda^2}{24}R_2-\frac{\lambda^3+4}{72}R_{11}.
\]
$T_{111}=T_{222}=T_{333}=2$, $T_{123}=\lambda$
および添字対称性を、$I,J$ の完全縮約へ代入して整理すると上の恒等式を得る。
これは全係数を独立記号とした有限多項式恒等式であり、
保存した検証器は左辺と右辺の全係数差が零であることを確認する。
ここで検算したのは固定入力の中心公式の Weyl symbol であり、
非可換中心性そのものをこの計算で証明したという意味ではない。
\end{proof}



\begin{theorem}[混合残余 potential の普遍四次係数]
\label{thm:nfm-mixed-fourth-potential}
五次 tensor の係数 $\beta=1/60$, $\gamma=1/240$ を用いると、
指定中心元に対して全 $n\ge3$ で
\[
 \boxed{\quad W_4=\Psi_\lambda(F)+\frac1{24}I(F)-\frac1{80}J(F)\quad}
\]
であり、同値に
\[
 W_4=-\frac S{80}F+\frac1{60}I+\frac1{240}J+\mathcal C_4(U).
\]
さらに Casimir 分解のスカラーは
\[
 a_h=h\left(1+\frac{S}{240}h^4\right)+O(h^7)
\]
である。従って先の表示における純中心項 $G_4$ は
$\mathcal C_4(U)$ に等しい。
\end{theorem}

\begin{proof}
局所共変式の補題によって $W_4=\Psi+dI+eJ$ と書く。
まず $U=0$ へ制限する。
$I(g)=J(g)=Sg$, $\Psi(g)=-Sg/48$ および
$\pi_5(g)=S\Pi_g/80$ を用いると
\[
 a_4=S(1/30-d-e)
\]
となる。$a_4$ は座標変数に依存しないので、この式は混合係数にもそのまま有効である。
前補題の恒等式を微分し、
$\pi_5=a_4\Pi_F+\Pi_{W_4}$ と比較する。
$S\,d_KF$, $d_KI$, $d_KJ$ は既に確認した独立な三つの共変式なので
\[
 d-\frac1{40}=\beta=\frac1{60},\qquad
 e+\frac1{60}=\gamma=\frac1{240}
\]
を得る。従って $d=1/24$, $e=-1/80$, $a_4=S/240$ である。
最後の純中心項の同定も、前補題の恒等式から従う。
\end{proof}

\begin{corollary}[四次の幾何的比較類]
\label{cor:nfm-fourth-jacobi-class}
$\tau_F$ を従来の八次 active Jacobi 代数の正規化 trace とする。
全変数の Jacobi 環で
\[
 \boxed{\quad
 [W_4-\tau_F(\Psi_\lambda(F))]
 =\frac1{60}[I(F)]+\frac1{40}[\tau_F(I(F))]
             +\frac S{80}[\tau_F(F)] .\quad}
\]
特に四次の幾何的比較には $\gamma$ と $a_4$ を別途評価する必要はなく、
純中心項の別のグラフ積分も不要である。
\end{corollary}

\begin{proof}
$J$ の縮約は $q_l=F_l$ を一つ含むので、
$J=\sum_l B_l F_l$ と書ける。各 $B_l$ は全座標の一次式である。
従って $[J]=0$ かつ $\tau_F(J)=0$。
また Euler の公式から全 Jacobi 環では $[F]=0$ である。
中心 symbol の恒等式に $\tau_F$ を適用して
\[
 \mathcal C_4-\tau_F(\Psi)
       =\frac S{80}\tau_F(F)+\frac1{40}\tau_F(I)
\]
を得る。これを四次 potential の公式へ代入すればよい。
\end{proof}

\begin{proposition}[Fermat active 切断上の混合公式と一般次元の非消滅]
\label{prop:nfm-fourth-uniform-witness}
$\lambda=0$ の上記全混合族では
\[
 [W_4-\tau_F(\Psi_0(F))]
 =\frac19[ACx+ABy+BCz+AH+BD+CE].
\]
従って各 $n\ge3$ の滑らかな全混合係数空間で、
この四次比較類が非零となる Zariski 開は非空である。
\end{proposition}

\begin{proof}
$R_{11}=ACx+ABy+BCz$, $K=AH+BD+CE$ と書く。
縮約式を $\lambda=0$ で展開すると
\[
 I=-\frac83(Az^2+Bx^2+Cy^2)+\frac43R_{11},\qquad
 J=-\frac23(AF_z+BF_x+CF_y).
\]
$F_x=x^2+Ay+Cz+D$ 等を使えば
\[
 I+\frac83(AF_z+BF_x+CF_y)
        =\frac{20}{3}R_{11}+\frac83K.
\]
八つの平方自由単項式を基底とする active Jacobi 代数では
$\tau_F(x)=\tau_F(y)=\tau_F(z)=0$ なので $\tau_F(I)=8K/3$ である。
前系を適用すると所望の $[R_{11}+K]/9$ を得る。

非消滅には、全次元に共通して
\[
 F_*=\frac13\left(x^3+y^3+z^3+u^3+v^3+\sum_j w_j^3\right)
                    +u xy+v yz
\]
を使える。$n=3$ では $w_j$ はなく、それ以外では $n-3$ 個を加える。
この三次は滑らかである。実際、全偏微分が零なら $w_j=0$。
$y=0$ なら残りも全て零なので、非零解があれば $y=1$ と規格化できる。
そのとき $x^2+u=u^2+x=0$ から $xu\in\{0,1\}$、
$z^2+v=v^2+z=0$ から $zv\in\{0,1\}$ となるが、
$F_y=1+xu+zv$ は零にならない。

三次多項式の空間上の線型汎関数 $\ell$ を
\[
 \ell(vx^2)=\ell(uz^2)=1,\qquad \ell(uvy)=-1
\]
とし、その他の三次単項式では零と定める。
これは全ての $z_i\partial_{z_j}F_*$ と全ての純中心三次を消す。
非自明な確認は
$vF_x=vx^2+uvy$ と $uF_z=uz^2+uvy$ の二つだけである。
従って全 Jacobi 環上の汎関数へ降り、
\[
 \ell\left(W_4-\tau_F(\Psi_0(F_*))\right)
      =\ell(uvy)/9=-1/9\ne0
\]
を得る。実際の比較類は滑らかな係数開上の Jacobi 束の正則切断なので、
非消滅は Zariski 開である。
ここでは純中心三次の像でさらに商を取った束の一定階数性を仮定していない。
\end{proof}

\begin{corollary}[Hessian 補正後にも残る標識付き比較の障害]
\label{cor:nfm-corrected-fourth-obstruction}
各 $n\ge3$ について、前命題の四次類が非零となる開の複素点で、
指定した中心量子化・PBW・formality の入力を満たすものを取る。
残余三次 $V(W_h)$ と有限補正三次 $V(H^\sharp)$ の間には、
同じ $h$ と特殊ファイバーの同定を保つ
$\C[h]/(h^5)$ 上の同型は存在しない。
固定入力の完備 MF 移送と相対 window 比較を適用する場合には、
量子残余圏と $\operatorname{Ku}(V(H^\sharp))$ の間にも、
同じ標識を保つ $\C[[h]]$-線型 dg Morita 同値は存在しない。
後者は $n=4$ を含む全 $n\ge3$ で成立する入力付きの主張である。
\end{corollary}

\begin{proof}
二つの三次は $h^4$ を法として同じである。
小拡大 $\C[h]/(h^5)\to\C[h]/(h^4)$ 上の持上げの差は
$h^4[W_4-\tau_F(\Psi)]$ である。
滑らかな三次の $H^0(T)=0$ により、特殊ファイバー上で恒等な同型は、
逐次の小拡大を通じて $h^4$ を法として恒等となる。
従ってそのような同型が存在すれば四次の差が Jacobi 環で零となり、矛盾する。

圏についても、命題\ref{prop:nf-hh}の $HH^1=0$ と
実際の幾何的 Kodaira--Spencer 写像
$(J_F)_3\hookrightarrow HH^2(\operatorname{Ku}(V(F)))$ の単射性を使う。
共通の $h^4$ 切断の持上げの差は、非零の四次類の像である。
特殊ファイバー上恒等な同値の曖昧さは $HH^1=0$ により逐次消えるため、
この差を同値で消せない。最後に、剰余で恒等かつ Artin 切断と両立する
固定入力の MF・window 同値を用いる。
四次元の幾何的変形の全射性は必要ない。
\end{proof}

この結論は、一般ファイバー上で標識を忘れた三次超曲面または
Kuznetsov 成分の非同値を意味しない。
また、Hessian 補正 $\Phi$ の全 O-降下と次数九の定理はそのまま有効であり、
ここで否定したのは $\Phi$ の三次と残余三次の全次数での標識付き一致である。

\subsubsection{四次まで適合する代数的 Poisson 射}

四次の差も座標不変に打ち消せる。
$P\in\det K$ を使い、$S_P(H),I_P(H),J_P(H)$ を
先の完全縮約の四つの $\varepsilon$ を $P$ に置き換えて定義する。
局所体積枠で $P=sP_0$ なら、それぞれ
$s^4S(H|_K),s^4I(H),s^4J(H)$ である。
$\tau=\tau_{H,K}$ と書き、
\begin{align}
 \mathcal E_P(H)
  &=\frac{I_P(H)}{60}+\frac{J_P(H)}{240}
     -\frac{S_P(H)}{80}\bigl(H-\tau(H)\bigr)
     +\frac{\tau(I_P(H))}{40},\label{eq:nfm-fourth-invariant-correction}\\
 H^{[4]}&=H^\sharp+\mathcal E_P(H),&
 P^{[4]}&=\left(1+\frac{S_P(H)}{240}\right)P
 \label{eq:nfm-fourth-corrected-pair}
\end{align}
と置く。

\begin{theorem}[全 O-同値へ降下する四次補正]
\label{thm:nfm-fourth-algebraic-correction}
$H,H|_K,H^{[4]},H^{[4]}|_K$ が滑らかで
$1+S_P(H)/240$ が可逆な非空開上で、
\[
 \Phi^{[4]}:(H,P)\longmapsto(H^{[4]},P^{[4]})
\]
は分解可能な三次 Poisson 対の代数的射を定める。
この射は任意の基底変更・方程式の生成元の変更・全射影同型と両立し、
既存の良い O-marked moduli の対応する開へ降下する。

指定中心元の trace--Hessian 像
$H=H^{\mathrm{TH}}_{F,h}$ と $P=hP_0$ に沿って、
固定 harmonic-angle 模型では全混合係数・全 $n\ge3$ に対して
\[
 H^{[4]}\equiv W_h\pmod{h^6},\qquad
 \iota_{dH^{[4]}}P^{[4]}\equiv\pi_h\pmod{h^7}
\]
が成立する。この有限式の射と、その全次数での残余比較は別の主張であり、
後者の収束・代数化・同値核の降下はここでは結論しない。
\end{theorem}

\begin{proof}
四つの縮約は $K$ の枠によらず、$\tau$ は任意底変更と両立する。
$(H,P)$ を同じ Poisson 対を表す $(cH,c^{-1}P)$ に替えると、
$S_P(H)$ は不変、$I_P(H),J_P(H)$ は $c$ 倍になる。
従って $\mathcal E_P(H)$ と $H^{[4]}$ は $c$ 倍、
$P^{[4]}$ は $c^{-1}$ 倍となる。全射影同型への共変性も同じ縮約から従う。
また、$H$ に純中心三次 $G(U)$ を加えても、
$S,I,J$ と active 臨界代数は変わらず、
$H-\tau(H)$ も変わらない。従って $\mathcal E_P$ はこの変更に不変である。

$P^{[4]}$ は $P$ の座標に依存しない単元倍なので、
出力は同じ rank 三 support を持つ Jacobian Poisson tensor を定める。
滑らかさと単元条件は開条件であり、$P=sP_0$ の $s=0$ では
$H^{[4]}=H$、単元は一となる。滑らかな三次と active 制限を出発点に取れば、
この形式 trait の一般点で必要な開条件を満たす。
既存の good 一ブロック開との交わりも非空である。
表示上の射を全同値へ降下する議論は
定理\ref{thm:nfm-algebraic-hessian-correction}と同じである。

jet の比較では $H^{\mathrm{TH}}_{F,h}-F$ が純中心三次であることを使う。
そのため
\[
 \mathcal E_{hP_0}(H^{\mathrm{TH}}_{F,h})
 =h^4\left(\frac I{60}+\frac J{240}
          -\frac S{80}(F-\tau_F(F))+\frac{\tau_F(I)}{40}\right)
 =h^4\bigl(W_4-\tau_F(\Psi)\bigr).
\]
$H^\sharp=F+h^2\mathscr H_F/24+h^4\tau_F(\Psi)$ と偶奇性から
最初の合同を得る。
さらに $P^{[4]}=h(1+Sh^4/240)P_0$ と
定理\ref{thm:nfm-mixed-fourth-potential}の $a_h$ の式を
$\pi_h=a_h\iota_{dW_h}P_0$ へ代入すれば二つ目の合同を得る。
\end{proof}

\begin{remark}[有限補正の二つの役割]
従来の $\Phi$ は量子側の楕円点曲線との厳密な一致と一般次数九を持つ。
新しい $\Phi^{[4]}$ についてここで証明したのは、全混合係数でのより高い残余 jet の一致である。
実際、純 active 三次では $\mathcal E_{hP_0}(g)=h^4S(g)g/120$ となるため、
新しい有限式に従来の点曲線との厳密な一致をそのまま移してはいない。
新しい射の一般次数も上の議論では決定していない。
\end{remark}


有限式
\[
 W_h^{[2]}=F+\frac{h^2}{24}\mathscr H_F,\qquad
 \pi_h^{[2]}=h\iota_{dW_h^{[2]}}P
\]
は全ての $h$ に対して Jacobi 恒等式と Casimir 条件を満たす。
しかし、固定 formality の実際の tensor との差は一般に五次から現れる。
四次 potential と新たな有限補正は上で決定したが、
三次の単元倍を許した全次数比較はなお別の課題である。

補正の座標共変性は内在的に書ける。
$K$ を rank 三の active 方向、$P\in\det K$ として
\[
 \mathscr H_{P}(F)
 =\bigl(\det\operatorname{Hess}_{K}F\bigr)(P\otimes P)
\]
と定める。局所的に $P=e_1\wedge e_2\wedge e_3$ と書けば
$\mathscr H_P(F)=\det(e_i e_jF)$ である。
従ってこれは補空間やその基底の選択に依存せず、
全線型座標変換に共変である。また
\[
 \mathscr H_{bP}(aF)=a^3b^2\mathscr H_P(F).
\]
特に $(F,P)\mapsto(aF,a^{-1}P)$ という方程式直線の局所自明化変更では、
$F+h^2\mathscr H_P(F)/24$ は $a$ 倍され、対応する tensor は変わらない。
従って有限 jet の補正は、方程式の代表や volume の選択から独立に定義される。

\begin{remark}[形式的な降下と残る代数化]
固定 formality と $h$ を保つ次数付き中心三次対の線型同型は、
正規化の一意性により形式 Poisson 対の同型を誘導する。
中心元を単元 $b$ 倍すると $W_h$ は $b$ 倍され、$\pi_h$ は変わらない。
従って構成は中心直線についても自然であり、
形式的な表示 groupoid 上では座標変更と cocycle に両立する。
この主張には formal Poisson tensor を単に候補として選ぶ曖昧さはない。

一方、ここで固定した formality 自体を変えた場合の比較、
形式級数の収束・代数化、$h$ の標識を忘れる操作、
および全 intrinsic O-marked 圏族への降下は別問題である。
Kontsevich\cite[Section~2.3, Conjecture~1 and Questions~1]{NFMKontsevichAlg}の
形式対応の収束・周期に関する問題とも区別する。
本節から有限 $h$ の一般点における代数的モジュライ射を宣言しない。
既存の全 O-同値比較を用いる場合も、形式変形とその標識を保つ同値が
表示 groupoid の同値へ移るという適用条件を確認してから降下する。
\end{remark}

\subsection{点曲線が決める split 族の厳密補正}

ここでは量子化の尺度 $s$ と Hesse 係数 $\lambda$ を独立に扱う。
点曲線の Hessian 表示は Mori--Smith\cite[Theorems~1.2, 7.2]{MoriSmith}
に既にある結果であり、新規性を主張する対象ではない。
本小節の用途は、その Hessian 表示を残余復元と結び付け、
補正後の Poisson 対の尺度と定義体を明記することである。

$\lambda\in\C$ を $1+\lambda^3\ne0$ なる定数とし、
$R=\C[s]_{(s)}$、$\widehat R=\C[[s]]$、$K=\C(s)$ と置く。
$S=x^3+y^3+z^3$、$P=xyz$、$g_\lambda=S/3+\lambda P$ と書く。
$D_R$ は
\[
 \begin{aligned}
 [y,z]&=s\left(x^2+\frac\lambda2(yz+zy)\right),\\
 [z,x]&=s\left(y^2+\frac\lambda2(zx+xz)\right),\\
 [x,y]&=s\left(z^2+\frac\lambda2(xy+yx)\right).
 \end{aligned}
\]
で定まる代数とする。$r=n-1\ge2$ とし、中心変数を加えて
$B_R=D_R[u_1,\ldots,u_r]$ と置く。
$G\in\C[u_1,\ldots,u_r]_3$ の偏微分の共通零点は原点だけと仮定する。
特に $G=\sum u_i^3/3$ を全ての $n$ に一様に使える。
\[
 \begin{aligned}
 a_{\lambda,s}&=1+\frac{s^2\lambda^2}{2}
                       +\frac{s^4\lambda}{6}-\frac{s^4\lambda^4}{48},\\
 b_{\lambda,s}&=\lambda-\frac{s^2\lambda^3}{6}-\frac{2s^2}{3}
                       +\frac{s^4\lambda^2}{6}-\frac{s^4\lambda^5}{48},\\
 \Omega_D&=a_{\lambda,s}S/3+b_{\lambda,s}\operatorname{sym}(P),
 \qquad w=\Omega_D+G.
 \end{aligned}
\]
これらは既存の Hesse 中心公式の二パラメータ表示である。
$r$ 個の中心変数を付加しても、各次数の有限自由性と中心性は保存される。

\begin{proposition}[点曲線の厳密な Hessian 補正]
\label{prop:nfm-hesse-point-correction}
$\mathscr H_\lambda=\det\operatorname{Hess}_{x,y,z}(g_\lambda)$ と置き、
\[
 \begin{aligned}
 A_{\lambda,s}&=1-\frac{s^2\lambda^2}{4},\\
 B_{\lambda,s}&=\lambda+\frac{s^2(4+\lambda^3)}{12},\\
 J_{\lambda,s}&=g_\lambda+\frac{s^2}{24}\mathscr H_\lambda
       =A_{\lambda,s}S/3+B_{\lambda,s}P,\\
 \nu(\lambda,s)&=\frac{B_{\lambda,s}}{A_{\lambda,s}}
       =\frac{\lambda+s^2(4+\lambda^3)/12}{1-s^2\lambda^2/4}
 \end{aligned}
\]
とする。一般点の active な点曲線は $V(J_{\lambda,s})$ である。
この曲線の $R$ 上の滑らかな閉包は $V(S+3\nu(\lambda,s)P)$ であり、
$s=0$ のファイバーは $V(g_\lambda)$ である。
ただし $s=0$ の可換代数自体の点スキームは $\mathbb P^2$ であり、
この閉包と混同しない。

Trace--Hessian 出力は $H^{\mathrm{TH}}=g_\lambda+G$ であり、
補正後の対
\[
 H^{\mathrm{corr}}=J_{\lambda,s}+G,
 \qquad
 \Pi^{\mathrm{corr}}
 =s\,\iota_{dH^{\mathrm{corr}}}
       (\partial_x\wedge\partial_y\wedge\partial_z)
\]
は $K$ 上の滑らかな三次超曲面と非零 Poisson 構造を定める。
その active な平面切断は、もとの代数の点曲線と埋込みも含めて一致する。
\end{proposition}

\begin{proof}
Sklyanin 係数は
$a=1-s\lambda/2$、$b=-1-s\lambda/2$、$c=-s$ である。
De Laet--Le Bruyn\cite[Section~2, preprint p.~3]{NFMDLL}の中心元
（本稿の基礎節の式\texttt{gre:eq:known-center}も参照）を
$C=xyz+yzx+zxy$、$D=yxz+xzy+zyx$ により
\[
 a(b^3-c^3)C+b(c^3-a^3)D+c(a^3-b^3)S
\]
と書き、関係式 $C-D=sS+(s\lambda/2)(C+D)$ を代入すると
$-9s\Omega_D$ を得る。
一般点は滑らかな楕円型の Sklyanin 代数であり、
剰余は三変数多項式環である。
両者の Hilbert 関数が一致するので、DVR 上の構造定理により全次数が自由である。
一般点の中心性は次数四の torsionfree 性により $R$ 上へ戻る。
この次数別自由性と多項式環への特殊化により、固定入力の
$\texttt{gcd:prop:automatic-as}$ の仮定も満たされる。
従って一般点の $B$ は両側 noetherian な Koszul AS 正則代数であり、
大域次元と Gorenstein parameter は $N=n+2$ である。
また $w\bmod s=g_\lambda+G\ne0$ なので $w$ は正則で、中心商も次数別に自由である。

点曲線の行列式は
\[
 -\frac{a^3+b^3+c^3}{3abc}
 =\frac{\lambda+s^2(4+\lambda^3)/12}{1-s^2\lambda^2/4}
\]
を与える。直接計算でも
\[
 \mathscr H_\lambda=(8+2\lambda^3)P-2\lambda^2S
\]
であり、表示した $J_{\lambda,s}$ を得る。
$A_{\lambda,s}$ と $1+\nu(\lambda,s)^3$ は $R$ の単元なので曲線は滑らかである。
点自己同型を与える切断は $(a:b:c)$ であり、$s=0$ では $(1:-1:0)$ に特殊化する。

active な Jacobi 環で $x^3=y^3=z^3=-\lambda P$ であるから
\[
 b_{\lambda,s}-\lambda a_{\lambda,s}
 =-\frac23s^2(1+\lambda^3),\qquad
 \nu_{s\Pi_{g_\lambda}}(\overline g_\lambda)
 =8s^2(1+\lambda^3)[P].
\]
正次数の active 元の乗法 trace は零であり、純中心成分は $G$ となる。
従って $-12$ 規格化の Trace--Hessian 出力は $g_\lambda+G$ である。
$H^{\mathrm{corr}}\bmod s=g_\lambda+G$ は滑らかなので、
$R$ 上の族と一般ファイバーは滑らかである。
Jacobian 形の三変数 Poisson tensor は全ての $u_i$ を Casimir とし、
$H^{\mathrm{corr}}$ も Casimir であるから、表示した対を得る。
\end{proof}

\begin{theorem}[split 族での残余圏と補正後三次の比較]
\label{thm:nfm-split-corrected-residual}
上の族について、固定入力論文Bの完備 MF 移送、相対 window、
固定対象の局所化および有限降下、coherent 残余圏の飽和性の入力を用いる。
特に最後の入力は $\texttt{crs:prop:residual-smooth}$ と
$\texttt{crs:thm:coherent}$ である。
このとき、完備残余復元の Hesse 係数は全次数で $\nu(\lambda,s)$ に等しい。
さらに、ある有限代数拡大 $K'/K$ 上で
\[
 (\mathcal K_{B/(w)},\rho)
 \simeq
 (\operatorname{Ku}(V(H^{\mathrm{corr}})),\rho_{\mathrm{corr}})
\]
という dg Morita 同値が存在する。
同値は $w$ と表示した $H^{\mathrm{corr}}$ に由来する指定三次周期を保存する。
これは任意の $n\ge3$ と任意の定数 $\lambda\in\C$、$1+\lambda^3\ne0$ に成立する。
\end{theorem}

\begin{proof}
固定入力の一パラメータ式へ単に記号を代入するのではなく、
その構成を本族について確認する。
まず三変数の $D_{\widehat R}$ と $\Omega_D$ に完備移送を適用する。
$g_\lambda$ は滑らかで、上で次数別自由性と中心性を確認したので、その仮定を満たす。
巡回置換と $(x,y,z)\mapsto(x,\omega y,\omega^2z)$ が生成する有限 Heisenberg 群は、
全ての $s$ に対して関係式と $\Omega_D$ を保存する。
ここでも harmonic-angle formality と、次数一生成元を固定する模型を使う。
補題\ref{lem:nfm-normalized-star-model}の一意性と線型共変性により、
この模型は自動的に Heisenberg 群と同変である。
従って固定入力の $\texttt{gre:prop:hesse-normalform}$ と同じ不変三次空間に入り、
active potential は
\[
 W_{\mathrm{act}}=a(s)S/3+b(s)P,
 \qquad a(0)=1,\quad b(0)=\lambda
\]
となる。中心多項式変数への自明な延長により、全 potential は
$W=W_{\mathrm{act}}+G$ である。
active 座標だけを $a(s)^{-1/3}$ 倍して得る Hesse 係数を
$\mu(s)=b(s)/a(s)$ と書く。
普遍二次 jet 計算により
\[
 a(s)=1-\frac{s^2\lambda^2}{4}+O(s^3),\qquad
 b(s)=\lambda+\frac{s^2(4+\lambda^3)}{12}+O(s^3),
\]
従って
\[
 \mu(s)=\lambda+\frac{1+\lambda^3}{3}s^2+O(s^3).
\]

$D_R/(\Omega_D)$ の相対 THCR 表示も同じ証明で得られる。
すなわち、滑らかな点曲線と平行移動切断を $R$ 上で取り、
その $\mathcal O(1)$ の twisted tensor 冪の切断環への写像を次数一から構成する。
各正次数の切断加群は有限自由であり、一般点で二次関係と $\Omega_D$ が消えるので、
torsionfree 性によりこの写像は $R$ 上で $D_R/(\Omega_D)$ を経由する。
また $\Omega_D\bmod s=g_\lambda$ は非零因子であるから、
次数別 Nakayama と分裂単射によって商の各次数も有限自由となる。点曲線と平行移動切断の剰余は
$(V(g_\lambda),(1:-1:0))$ であり、剰余の商写像は通常の斉次座標環の同型である。
各次数で同じ階数の有限自由加群の間の写像なので Nakayama により同型となる。
相対 ampleness と三変数の指数零の MF 比較は固定入力の
$\texttt{gre:prop:elliptic-quotient}$、$\texttt{gre:lem:elliptic-mf}$ の証明通りである。
ここでは Fermat という条件を使わず、$V(g_\lambda)$ の滑らかさだけを使う。
従って $\C((s))$ 上で $E_\mu$ と $E_\nu$ の導来同値を得る。
両者は $(1:-1:0)$ を原点に持つ楕円曲線なので、
Antieau--Krashen--Ward\cite[Theorem~2.2]{NFMAKW}により同型となり、
$j(\mu)=j(\nu)$ である。

この等式から分枝を一意に決める。
$j(z)$ は $z=\lambda$ で正則な非定数有理関数なので、ある整数 $e\ge1$ と
形式的可逆変数 $\psi(z)=z+O(z^2)$、非零定数 $c$ により
\[
 j(\lambda+z)-j(\lambda)=c\,\psi(z)^e
\]
と書ける。$\mu-\lambda$ と $\nu-\lambda$ はともに
$(1+\lambda^3)s^2/3+O(s^3)$ で始まり、その係数は非零である。
$j$ の等式から二つの $\psi$ の値の比は $e$ 乗根となり、
その剰余は一だから比は一である。
$\psi$ の形式逆により $\mu=\nu$ を得る。
この証明は $j'(\lambda)=0$ の場合も含む。

従って全 potential は $a(s)g_\nu+G$ である。
active 座標だけを $(A_{\lambda,s}/a(s))^{1/3}$ 倍すれば、
厳密に $H^{\mathrm{corr}}=A_{\lambda,s}g_\nu+G$ へ移る。
この根は剰余が一であり、底の $s$、中心変数、周期変数 $t$ を固定する。
従って内部 shift と指定三次周期を保つ完備 MF 比較が得られる。

$H^{\mathrm{corr}}$ 自体は $K$ 上の多項式である。
固定入力の一般次元の固定対象への制限、局所化、生成性および有限降下を
この既に指定した代数的標的へ適用すると、ある有限 $K'/K$ 上で
まず線束生成残余部分について同値を得る。
ここで用いる局所化と有限降下は $\texttt{rog:lem:localization}$ と
$\texttt{rog:lem:descent}$ である。
全 coherent 残余圏への移行には、さらに上記の残余 smoothness と
coherent 飽和性の入力を用いる。これにより表示した $\mathcal K$ 全体の同値となる。
完備化上の形式座標変換そのものが $K'$ 上に代数化されることは使わない。
また、特殊点を含む $R$ 上の代数的相対圏同値はここでは結論しない。
\end{proof}

\begin{remark}[尺度と既存の一パラメータ例]
$A=A_{\lambda,s}$ と書く。
$A$ の三乗根を加え、active 座標に $x,y,z\mapsto A^{-1/3}(x,y,z)$ を施すと、
補正後の対は
\[
 \left(g_\nu+G,\ sA\,\Pi_{g_\nu}\right)
\]
へ移る。三次式を規格化するだけで tensor を $s\Pi_{g_\nu}$ と書くのは正しくない。
三重ベクトルは逆 determinant によって $A$ 倍されるためである。
補正前の表示 $(H^{\mathrm{corr}},\Pi^{\mathrm{corr}})$ は既に $K$ 上にあり、
三乗根は規格化した表示との同型のためにだけ必要となる。

前節の $\lambda=s=h$ の族では
$A=1-h^4/4$、$\nu=\lambda_{\mathrm{res}}(h)$ となる。
従って補正後の代数的三次は、前節で独立に決定した残余三次と
有限三乗根拡大後に厳密な座標変換で一致する。
この場合の残余比較は、前節の固定入力の exact formula からも直接従う。

上の $\Pi^{\mathrm{corr}}$ は、点曲線と Hessian 共変式によって選んだ
代数的 Poisson tensor である。
固定した Kontsevich 比較の formal Poisson tensor の全次数と同一だとは証明していない。
後者の尺度や gauge の全高次比較、一般混合係数での残余比較、
全 O-同値に対する残余同値の降下は別の課題として残る。
\end{remark}
